% Figura (simulación): electrete cilíndrico uniformemente polarizado (problema 4.4),
% líneas de P, E y D.  Datos: scripts/lineas_campo.py -> datos/cilindro-polarizado-{E,D}.tex
\begin{figure}[t]
\centering
\pgfplotsset{cil/.style={marco, width=5.55cm, xmin=-3.2, xmax=3.2, ymin=-1.9, ymax=1.9}}
\begin{subfigure}[b]{0.325\textwidth}\centering
\begin{tikzpicture}
\begin{axis}[cil]
  \path (rel axis cs:0,0) rectangle (rel axis cs:1,1);
  \fill[dielec] (axis cs:-1.5,-0.5) rectangle (axis cs:1.5,0.5);
  \pgfplotsinvokeforeach{-0.4,-0.2,0,0.2,0.4}{
    \addplot[draw=campoP, line width=0.6pt, flechaen=0.55] coordinates {(-1.5,#1) (1.5,#1)};}
  \draw[texto, line width=0.6pt] (axis cs:-1.5,-0.5) rectangle (axis cs:1.5,0.5);
\end{axis}
\end{tikzpicture}
\caption{$\vP(\vr)$}
\end{subfigure}\hfill
\begin{subfigure}[b]{0.325\textwidth}\centering
\begin{tikzpicture}
\begin{axis}[cil]
  \fill[dielec] (axis cs:-1.5,-0.5) rectangle (axis cs:1.5,0.5);
  %% Generado por scripts/lineas_campo.py -- no editar a mano
\addplot[equipot] coordinates {(-1.500,-0.395) (-1.516,-0.364) (-1.531,-0.332) (-1.545,-0.300) (-1.557,-0.267) (-1.568,-0.234) (-1.577,-0.200) (-1.585,-0.166) (-1.591,-0.131) (-1.596,-0.097) (-1.599,-0.062) (-1.601,-0.027) (-1.601,0.008) (-1.600,0.043) (-1.598,0.078) (-1.594,0.113) (-1.588,0.147) (-1.581,0.182) (-1.573,0.216) (-1.563,0.249) (-1.552,0.282) (-1.539,0.315) (-1.524,0.347) (-1.508,0.378) (-1.494,0.380) (-1.477,0.349) (-1.463,0.317) (-1.451,0.284) (-1.439,0.251) (-1.430,0.217) (-1.422,0.183) (-1.415,0.149) (-1.410,0.114) (-1.406,0.079) (-1.403,0.044) (-1.402,0.009) (-1.402,-0.026) (-1.404,-0.061) (-1.407,-0.095) (-1.412,-0.130) (-1.418,-0.165) (-1.425,-0.199) (-1.434,-0.233) (-1.444,-0.266) (-1.456,-0.299) (-1.470,-0.332) (-1.484,-0.363) (-1.500,-0.395)};
\addplot[equipot] coordinates {(-1.500,-0.482) (-1.526,-0.459) (-1.550,-0.434) (-1.572,-0.406) (-1.592,-0.378) (-1.610,-0.347) (-1.627,-0.316) (-1.641,-0.284) (-1.654,-0.252) (-1.665,-0.218) (-1.674,-0.185) (-1.682,-0.150) (-1.688,-0.116) (-1.692,-0.081) (-1.695,-0.046) (-1.697,-0.011) (-1.696,0.024) (-1.694,0.059) (-1.691,0.094) (-1.686,0.129) (-1.679,0.163) (-1.671,0.197) (-1.661,0.231) (-1.649,0.264) (-1.636,0.297) (-1.621,0.328) (-1.604,0.359) (-1.585,0.389) (-1.564,0.417) (-1.541,0.444) (-1.516,0.468) (-1.490,0.474) (-1.465,0.449) (-1.442,0.422) (-1.422,0.394) (-1.403,0.364) (-1.386,0.333) (-1.372,0.301) (-1.358,0.269) (-1.347,0.235) (-1.337,0.202) (-1.329,0.168) (-1.323,0.133) (-1.317,0.098) (-1.314,0.063) (-1.312,0.028) (-1.312,-0.007) (-1.313,-0.042) (-1.315,-0.077) (-1.319,-0.112) (-1.325,-0.146) (-1.332,-0.181) (-1.341,-0.215) (-1.351,-0.248) (-1.363,-0.281) (-1.377,-0.313) (-1.392,-0.345) (-1.410,-0.375) (-1.429,-0.405) (-1.451,-0.432) (-1.474,-0.458) (-1.500,-0.482)};
\addplot[equipot] coordinates {(-1.548,-0.528) (-1.580,-0.514) (-1.610,-0.495) (-1.637,-0.472) (-1.662,-0.448) (-1.685,-0.421) (-1.707,-0.393) (-1.726,-0.363) (-1.743,-0.333) (-1.759,-0.301) (-1.773,-0.269) (-1.785,-0.236) (-1.796,-0.202) (-1.804,-0.168) (-1.812,-0.133) (-1.817,-0.098) (-1.821,-0.063) (-1.823,-0.028) (-1.823,0.007) (-1.822,0.043) (-1.819,0.078) (-1.815,0.113) (-1.809,0.148) (-1.801,0.182) (-1.791,0.216) (-1.780,0.250) (-1.767,0.282) (-1.753,0.315) (-1.736,0.346) (-1.718,0.376) (-1.698,0.405) (-1.676,0.432) (-1.652,0.458) (-1.626,0.482) (-1.597,0.503) (-1.567,0.520) (-1.534,0.533) (-1.499,0.537) (-1.464,0.531) (-1.432,0.517) (-1.402,0.498) (-1.375,0.475) (-1.350,0.450) (-1.327,0.423) (-1.307,0.395) (-1.288,0.365) (-1.271,0.334) (-1.256,0.302) (-1.243,0.269) (-1.231,0.236) (-1.221,0.202) (-1.213,0.168) (-1.206,0.133) (-1.201,0.098) (-1.197,0.063) (-1.196,0.028) (-1.195,-0.007) (-1.196,-0.043) (-1.199,-0.078) (-1.203,-0.113) (-1.209,-0.148) (-1.216,-0.182) (-1.225,-0.216) (-1.236,-0.250) (-1.248,-0.283) (-1.262,-0.316) (-1.278,-0.347) (-1.295,-0.378) (-1.315,-0.407) (-1.337,-0.435) (-1.360,-0.461) (-1.386,-0.485) (-1.414,-0.506) (-1.445,-0.523) (-1.478,-0.535) (-1.514,-0.536) (-1.548,-0.528)};
\addplot[equipot] coordinates {(-1.548,-0.656) (-1.583,-0.651) (-1.617,-0.642) (-1.650,-0.630) (-1.682,-0.616) (-1.712,-0.599) (-1.742,-0.579) (-1.770,-0.558) (-1.796,-0.535) (-1.821,-0.510) (-1.845,-0.484) (-1.867,-0.457) (-1.887,-0.428) (-1.906,-0.399) (-1.924,-0.368) (-1.940,-0.337) (-1.954,-0.305) (-1.967,-0.272) (-1.978,-0.239) (-1.988,-0.205) (-1.996,-0.171) (-2.003,-0.137) (-2.008,-0.102) (-2.011,-0.067) (-2.014,-0.032) (-2.014,0.003) (-2.013,0.039) (-2.011,0.074) (-2.007,0.108) (-2.001,0.143) (-1.994,0.178) (-1.986,0.212) (-1.976,0.245) (-1.964,0.279) (-1.951,0.311) (-1.937,0.343) (-1.920,0.374) (-1.903,0.405) (-1.883,0.434) (-1.863,0.462) (-1.840,0.490) (-1.817,0.515) (-1.791,0.540) (-1.764,0.562) (-1.736,0.583) (-1.707,0.602) (-1.676,0.619) (-1.643,0.633) (-1.610,0.644) (-1.576,0.652) (-1.541,0.657) (-1.506,0.658) (-1.471,0.656) (-1.436,0.650) (-1.403,0.640) (-1.370,0.628) (-1.338,0.612) (-1.308,0.594) (-1.279,0.574) (-1.252,0.552) (-1.227,0.527) (-1.203,0.501) (-1.181,0.474) (-1.160,0.446) (-1.141,0.416) (-1.124,0.386) (-1.108,0.354) (-1.094,0.322) (-1.081,0.289) (-1.070,0.256) (-1.060,0.222) (-1.052,0.188) (-1.045,0.154) (-1.039,0.119) (-1.035,0.084) (-1.032,0.049) (-1.031,0.014) (-1.031,-0.021) (-1.033,-0.056) (-1.036,-0.091) (-1.040,-0.126) (-1.046,-0.161) (-1.053,-0.195) (-1.062,-0.229) (-1.072,-0.263) (-1.084,-0.296) (-1.097,-0.328) (-1.111,-0.360) (-1.127,-0.392) (-1.145,-0.422) (-1.164,-0.451) (-1.185,-0.480) (-1.208,-0.507) (-1.232,-0.532) (-1.257,-0.556) (-1.285,-0.578) (-1.314,-0.598) (-1.344,-0.616) (-1.376,-0.631) (-1.409,-0.642) (-1.443,-0.651) (-1.478,-0.656) (-1.513,-0.658) (-1.548,-0.656)};
\addplot[equipot] coordinates {(-1.660,-0.922) (-1.695,-0.915) (-1.729,-0.907) (-1.763,-0.898) (-1.796,-0.887) (-1.829,-0.874) (-1.861,-0.860) (-1.893,-0.844) (-1.924,-0.827) (-1.954,-0.809) (-1.983,-0.789) (-2.011,-0.769) (-2.039,-0.747) (-2.065,-0.724) (-2.091,-0.699) (-2.115,-0.674) (-2.139,-0.648) (-2.162,-0.621) (-2.183,-0.593) (-2.204,-0.565) (-2.223,-0.535) (-2.241,-0.505) (-2.258,-0.474) (-2.274,-0.443) (-2.289,-0.411) (-2.303,-0.379) (-2.315,-0.346) (-2.327,-0.313) (-2.337,-0.279) (-2.346,-0.245) (-2.354,-0.211) (-2.360,-0.176) (-2.366,-0.141) (-2.370,-0.106) (-2.373,-0.071) (-2.375,-0.036) (-2.376,-0.001) (-2.375,0.034) (-2.373,0.069) (-2.370,0.104) (-2.366,0.139) (-2.361,0.174) (-2.354,0.209) (-2.346,0.243) (-2.337,0.277) (-2.327,0.311) (-2.316,0.344) (-2.304,0.377) (-2.290,0.409) (-2.275,0.441) (-2.259,0.473) (-2.242,0.503) (-2.224,0.534) (-2.205,0.563) (-2.184,0.592) (-2.163,0.620) (-2.140,0.647) (-2.117,0.673) (-2.092,0.698) (-2.067,0.722) (-2.040,0.745) (-2.013,0.767) (-1.985,0.788) (-1.955,0.808) (-1.925,0.826) (-1.895,0.843) (-1.863,0.859) (-1.831,0.873) (-1.798,0.886) (-1.765,0.897) (-1.731,0.907) (-1.697,0.915) (-1.662,0.921) (-1.627,0.926) (-1.592,0.929) (-1.557,0.930) (-1.522,0.930) (-1.487,0.927) (-1.452,0.923) (-1.417,0.917) (-1.383,0.909) (-1.349,0.900) (-1.316,0.889) (-1.283,0.876) (-1.251,0.861) (-1.219,0.846) (-1.189,0.828) (-1.159,0.809) (-1.130,0.789) (-1.103,0.767) (-1.076,0.744) (-1.050,0.720) (-1.026,0.695) (-1.002,0.669) (-0.980,0.642) (-0.959,0.613) (-0.939,0.584) (-0.920,0.555) (-0.902,0.524) (-0.886,0.493) (-0.870,0.462) (-0.856,0.430) (-0.843,0.397) (-0.831,0.364) (-0.820,0.330) (-0.810,0.297) (-0.801,0.263) (-0.794,0.228) (-0.787,0.194) (-0.782,0.159) (-0.778,0.124) (-0.774,0.089) (-0.772,0.054) (-0.771,0.019) (-0.771,-0.017) (-0.772,-0.052) (-0.774,-0.087) (-0.777,-0.122) (-0.782,-0.157) (-0.787,-0.192) (-0.793,-0.226) (-0.801,-0.261) (-0.809,-0.295) (-0.819,-0.329) (-0.830,-0.362) (-0.842,-0.395) (-0.855,-0.428) (-0.869,-0.460) (-0.885,-0.492) (-0.901,-0.523) (-0.919,-0.553) (-0.938,-0.583) (-0.958,-0.612) (-0.979,-0.640) (-1.001,-0.667) (-1.024,-0.693) (-1.049,-0.719) (-1.075,-0.743) (-1.101,-0.766) (-1.129,-0.788) (-1.158,-0.808) (-1.187,-0.827) (-1.218,-0.845) (-1.249,-0.861) (-1.281,-0.875) (-1.314,-0.888) (-1.347,-0.899) (-1.381,-0.909) (-1.415,-0.917) (-1.450,-0.923) (-1.485,-0.927) (-1.520,-0.929) (-1.555,-0.930) (-1.590,-0.929) (-1.625,-0.926) (-1.660,-0.922)};
\addplot[equipot] coordinates {(-3.200,1.227) (-3.178,1.254) (-3.155,1.281) (-3.131,1.307) (-3.107,1.333) (-3.083,1.358) (-3.058,1.382) (-3.033,1.407) (-3.007,1.430) (-2.981,1.454) (-2.954,1.476) (-2.927,1.499) (-2.899,1.520) (-2.871,1.541) (-2.843,1.562) (-2.814,1.582) (-2.785,1.601) (-2.755,1.620) (-2.725,1.639) (-2.695,1.656) (-2.665,1.673) (-2.634,1.690) (-2.602,1.706) (-2.571,1.721) (-2.539,1.735) (-2.507,1.749) (-2.474,1.762) (-2.441,1.775) (-2.408,1.787) (-2.375,1.798) (-2.342,1.808) (-2.308,1.818) (-2.274,1.827) (-2.240,1.835) (-2.206,1.843) (-2.171,1.850) (-2.137,1.856) (-2.102,1.861) (-2.067,1.866) (-2.033,1.869) (-1.998,1.872) (-1.963,1.875) (-1.928,1.876) (-1.893,1.877) (-1.857,1.877) (-1.822,1.876) (-1.787,1.874) (-1.752,1.872) (-1.718,1.868) (-1.683,1.864) (-1.648,1.859) (-1.613,1.854) (-1.579,1.847) (-1.545,1.840) (-1.511,1.832) (-1.477,1.823) (-1.443,1.813) (-1.409,1.803) (-1.376,1.791) (-1.343,1.779) (-1.311,1.766) (-1.278,1.753) (-1.247,1.738) (-1.215,1.723) (-1.184,1.707) (-1.153,1.690) (-1.122,1.673) (-1.093,1.655) (-1.063,1.636) (-1.034,1.616) (-1.005,1.596) (-0.977,1.575) (-0.950,1.553) (-0.923,1.531) (-0.896,1.508) (-0.870,1.484) (-0.845,1.460) (-0.820,1.435) (-0.796,1.410) (-0.772,1.384) (-0.749,1.358) (-0.727,1.331) (-0.705,1.303) (-0.684,1.275) (-0.663,1.247) (-0.643,1.218) (-0.624,1.189) (-0.605,1.159) (-0.588,1.129) (-0.570,1.098) (-0.553,1.067) (-0.537,1.036) (-0.522,1.005) (-0.507,0.973) (-0.493,0.941) (-0.479,0.909) (-0.466,0.876) (-0.454,0.843) (-0.442,0.810) (-0.431,0.777) (-0.420,0.744) (-0.410,0.710) (-0.401,0.676) (-0.392,0.642) (-0.383,0.608) (-0.375,0.574) (-0.368,0.540) (-0.361,0.506) (-0.355,0.471) (-0.349,0.436) (-0.344,0.402) (-0.339,0.367) (-0.334,0.332) (-0.330,0.297) (-0.327,0.263) (-0.324,0.228) (-0.321,0.193) (-0.319,0.158) (-0.317,0.123) (-0.316,0.088) (-0.315,0.053) (-0.315,0.018) (-0.315,-0.018) (-0.315,-0.053) (-0.316,-0.088) (-0.317,-0.123) (-0.319,-0.158) (-0.321,-0.193) (-0.324,-0.228) (-0.327,-0.263) (-0.330,-0.297) (-0.334,-0.332) (-0.339,-0.367) (-0.344,-0.402) (-0.349,-0.436) (-0.355,-0.471) (-0.361,-0.506) (-0.368,-0.540) (-0.375,-0.574) (-0.383,-0.608) (-0.392,-0.642) (-0.401,-0.676) (-0.410,-0.710) (-0.420,-0.744) (-0.431,-0.777) (-0.442,-0.810) (-0.454,-0.843) (-0.466,-0.876) (-0.479,-0.909) (-0.493,-0.941) (-0.507,-0.973) (-0.522,-1.005) (-0.537,-1.036) (-0.553,-1.067) (-0.570,-1.098) (-0.588,-1.129) (-0.605,-1.159) (-0.624,-1.189) (-0.643,-1.218) (-0.663,-1.247) (-0.684,-1.275) (-0.705,-1.303) (-0.727,-1.331) (-0.749,-1.358) (-0.772,-1.384) (-0.796,-1.410) (-0.820,-1.435) (-0.845,-1.460) (-0.870,-1.484) (-0.896,-1.508) (-0.923,-1.531) (-0.950,-1.553) (-0.977,-1.575) (-1.005,-1.596) (-1.034,-1.616) (-1.063,-1.636) (-1.093,-1.655) (-1.122,-1.673) (-1.153,-1.690) (-1.184,-1.707) (-1.215,-1.723) (-1.247,-1.738) (-1.278,-1.753) (-1.311,-1.766) (-1.343,-1.779) (-1.376,-1.791) (-1.409,-1.803) (-1.443,-1.813) (-1.477,-1.823) (-1.511,-1.832) (-1.545,-1.840) (-1.579,-1.847) (-1.613,-1.854) (-1.648,-1.859) (-1.683,-1.864) (-1.718,-1.868) (-1.752,-1.872) (-1.787,-1.874) (-1.822,-1.876) (-1.857,-1.877) (-1.893,-1.877) (-1.928,-1.876) (-1.963,-1.875) (-1.998,-1.872) (-2.033,-1.869) (-2.067,-1.866) (-2.102,-1.861) (-2.137,-1.856) (-2.171,-1.850) (-2.206,-1.843) (-2.240,-1.835) (-2.274,-1.827) (-2.308,-1.818) (-2.342,-1.808) (-2.375,-1.798) (-2.408,-1.787) (-2.441,-1.775) (-2.474,-1.762) (-2.507,-1.749) (-2.539,-1.735) (-2.571,-1.721) (-2.602,-1.706) (-2.634,-1.690) (-2.665,-1.673) (-2.695,-1.656) (-2.725,-1.639) (-2.755,-1.620) (-2.785,-1.601) (-2.814,-1.582) (-2.843,-1.562) (-2.871,-1.541) (-2.899,-1.520) (-2.927,-1.499) (-2.954,-1.476) (-2.981,-1.454) (-3.007,-1.430) (-3.033,-1.407) (-3.058,-1.382) (-3.083,-1.358) (-3.107,-1.333) (-3.131,-1.307) (-3.155,-1.281) (-3.178,-1.254) (-3.200,-1.227)};
\addplot[equipot] coordinates {(3.200,1.227) (3.178,1.254) (3.155,1.281) (3.131,1.307) (3.107,1.333) (3.083,1.358) (3.058,1.382) (3.033,1.407) (3.007,1.430) (2.981,1.454) (2.954,1.476) (2.927,1.499) (2.899,1.520) (2.871,1.541) (2.843,1.562) (2.814,1.582) (2.785,1.601) (2.755,1.620) (2.725,1.639) (2.695,1.656) (2.665,1.673) (2.634,1.690) (2.602,1.706) (2.571,1.721) (2.539,1.735) (2.507,1.749) (2.474,1.762) (2.441,1.775) (2.408,1.787) (2.375,1.798) (2.342,1.808) (2.308,1.818) (2.274,1.827) (2.240,1.835) (2.206,1.843) (2.171,1.850) (2.137,1.856) (2.102,1.861) (2.067,1.866) (2.033,1.869) (1.998,1.872) (1.963,1.875) (1.928,1.876) (1.893,1.877) (1.857,1.877) (1.822,1.876) (1.787,1.874) (1.752,1.872) (1.718,1.868) (1.683,1.864) (1.648,1.859) (1.613,1.854) (1.579,1.847) (1.545,1.840) (1.511,1.832) (1.477,1.823) (1.443,1.813) (1.409,1.803) (1.376,1.791) (1.343,1.779) (1.311,1.766) (1.278,1.753) (1.247,1.738) (1.215,1.723) (1.184,1.707) (1.153,1.690) (1.122,1.673) (1.093,1.655) (1.063,1.636) (1.034,1.616) (1.005,1.596) (0.977,1.575) (0.950,1.553) (0.923,1.531) (0.896,1.508) (0.870,1.484) (0.845,1.460) (0.820,1.435) (0.796,1.410) (0.772,1.384) (0.749,1.358) (0.727,1.331) (0.705,1.303) (0.684,1.275) (0.663,1.247) (0.643,1.218) (0.624,1.189) (0.605,1.159) (0.588,1.129) (0.570,1.098) (0.553,1.067) (0.537,1.036) (0.522,1.005) (0.507,0.973) (0.493,0.941) (0.479,0.909) (0.466,0.876) (0.454,0.843) (0.442,0.810) (0.431,0.777) (0.420,0.744) (0.410,0.710) (0.401,0.676) (0.392,0.642) (0.383,0.608) (0.375,0.574) (0.368,0.540) (0.361,0.506) (0.355,0.471) (0.349,0.436) (0.344,0.402) (0.339,0.367) (0.334,0.332) (0.330,0.297) (0.327,0.263) (0.324,0.228) (0.321,0.193) (0.319,0.158) (0.317,0.123) (0.316,0.088) (0.315,0.053) (0.315,0.018) (0.315,-0.018) (0.315,-0.053) (0.316,-0.088) (0.317,-0.123) (0.319,-0.158) (0.321,-0.193) (0.324,-0.228) (0.327,-0.263) (0.330,-0.297) (0.334,-0.332) (0.339,-0.367) (0.344,-0.402) (0.349,-0.436) (0.355,-0.471) (0.361,-0.506) (0.368,-0.540) (0.375,-0.574) (0.383,-0.608) (0.392,-0.642) (0.401,-0.676) (0.410,-0.710) (0.420,-0.744) (0.431,-0.777) (0.442,-0.810) (0.454,-0.843) (0.466,-0.876) (0.479,-0.909) (0.493,-0.941) (0.507,-0.973) (0.522,-1.005) (0.537,-1.036) (0.553,-1.067) (0.570,-1.098) (0.588,-1.129) (0.605,-1.159) (0.624,-1.189) (0.643,-1.218) (0.663,-1.247) (0.684,-1.275) (0.705,-1.303) (0.727,-1.331) (0.749,-1.358) (0.772,-1.384) (0.796,-1.410) (0.820,-1.435) (0.845,-1.460) (0.870,-1.484) (0.896,-1.508) (0.923,-1.531) (0.950,-1.553) (0.977,-1.575) (1.005,-1.596) (1.034,-1.616) (1.063,-1.636) (1.093,-1.655) (1.122,-1.673) (1.153,-1.690) (1.184,-1.707) (1.215,-1.723) (1.247,-1.738) (1.278,-1.753) (1.311,-1.766) (1.343,-1.779) (1.376,-1.791) (1.409,-1.803) (1.443,-1.813) (1.477,-1.823) (1.511,-1.832) (1.545,-1.840) (1.579,-1.847) (1.613,-1.854) (1.648,-1.859) (1.683,-1.864) (1.718,-1.868) (1.752,-1.872) (1.787,-1.874) (1.822,-1.876) (1.857,-1.877) (1.893,-1.877) (1.928,-1.876) (1.963,-1.875) (1.998,-1.872) (2.033,-1.869) (2.067,-1.866) (2.102,-1.861) (2.137,-1.856) (2.171,-1.850) (2.206,-1.843) (2.240,-1.835) (2.274,-1.827) (2.308,-1.818) (2.342,-1.808) (2.375,-1.798) (2.408,-1.787) (2.441,-1.775) (2.474,-1.762) (2.507,-1.749) (2.539,-1.735) (2.571,-1.721) (2.602,-1.706) (2.634,-1.690) (2.665,-1.673) (2.695,-1.656) (2.725,-1.639) (2.755,-1.620) (2.785,-1.601) (2.814,-1.582) (2.843,-1.562) (2.871,-1.541) (2.899,-1.520) (2.927,-1.499) (2.954,-1.476) (2.981,-1.454) (3.007,-1.430) (3.033,-1.407) (3.058,-1.382) (3.083,-1.358) (3.107,-1.333) (3.131,-1.307) (3.155,-1.281) (3.178,-1.254) (3.200,-1.227)};
\addplot[equipot] coordinates {(1.452,-0.923) (1.487,-0.927) (1.522,-0.930) (1.557,-0.930) (1.592,-0.929) (1.627,-0.926) (1.662,-0.921) (1.697,-0.915) (1.731,-0.907) (1.765,-0.897) (1.798,-0.886) (1.831,-0.873) (1.863,-0.859) (1.894,-0.844) (1.925,-0.826) (1.955,-0.808) (1.984,-0.788) (2.013,-0.768) (2.040,-0.745) (2.067,-0.722) (2.092,-0.698) (2.117,-0.673) (2.140,-0.647) (2.163,-0.620) (2.184,-0.592) (2.205,-0.563) (2.224,-0.534) (2.242,-0.504) (2.259,-0.473) (2.275,-0.441) (2.290,-0.410) (2.304,-0.377) (2.316,-0.344) (2.327,-0.311) (2.337,-0.277) (2.346,-0.243) (2.354,-0.209) (2.361,-0.174) (2.366,-0.139) (2.370,-0.105) (2.373,-0.070) (2.375,-0.034) (2.376,0.001) (2.375,0.036) (2.373,0.071) (2.370,0.106) (2.366,0.141) (2.360,0.176) (2.354,0.210) (2.346,0.245) (2.337,0.279) (2.327,0.312) (2.315,0.346) (2.303,0.379) (2.289,0.411) (2.274,0.443) (2.258,0.474) (2.241,0.505) (2.223,0.535) (2.204,0.565) (2.183,0.593) (2.162,0.621) (2.139,0.648) (2.116,0.674) (2.091,0.699) (2.065,0.723) (2.039,0.747) (2.011,0.769) (1.983,0.789) (1.954,0.809) (1.924,0.827) (1.893,0.844) (1.861,0.860) (1.829,0.874) (1.796,0.887) (1.763,0.898) (1.729,0.907) (1.695,0.915) (1.660,0.922) (1.625,0.926) (1.590,0.929) (1.555,0.930) (1.520,0.929) (1.485,0.927) (1.450,0.923) (1.415,0.917) (1.381,0.909) (1.347,0.899) (1.314,0.888) (1.281,0.875) (1.249,0.861) (1.218,0.845) (1.187,0.827) (1.158,0.808) (1.129,0.788) (1.101,0.766) (1.075,0.743) (1.049,0.719) (1.025,0.694) (1.001,0.667) (0.979,0.640) (0.958,0.612) (0.938,0.583) (0.919,0.553) (0.901,0.523) (0.885,0.492) (0.869,0.460) (0.855,0.428) (0.842,0.395) (0.830,0.362) (0.819,0.329) (0.809,0.295) (0.801,0.261) (0.793,0.226) (0.787,0.192) (0.782,0.157) (0.777,0.122) (0.774,0.087) (0.772,0.052) (0.771,0.017) (0.771,-0.018) (0.772,-0.054) (0.774,-0.089) (0.778,-0.124) (0.782,-0.159) (0.787,-0.193) (0.794,-0.228) (0.801,-0.262) (0.810,-0.296) (0.820,-0.330) (0.831,-0.364) (0.843,-0.397) (0.856,-0.429) (0.870,-0.462) (0.885,-0.493) (0.902,-0.524) (0.920,-0.555) (0.939,-0.584) (0.959,-0.613) (0.980,-0.641) (1.002,-0.669) (1.026,-0.695) (1.050,-0.720) (1.076,-0.744) (1.103,-0.767) (1.130,-0.789) (1.159,-0.809) (1.189,-0.828) (1.219,-0.845) (1.251,-0.861) (1.283,-0.876) (1.315,-0.889) (1.349,-0.900) (1.383,-0.909) (1.417,-0.917) (1.452,-0.923)};
\addplot[equipot] coordinates {(1.468,-0.655) (1.503,-0.658) (1.538,-0.657) (1.573,-0.653) (1.607,-0.645) (1.640,-0.634) (1.673,-0.620) (1.704,-0.604) (1.733,-0.585) (1.762,-0.564) (1.789,-0.542) (1.814,-0.518) (1.838,-0.492) (1.861,-0.465) (1.882,-0.437) (1.901,-0.407) (1.919,-0.377) (1.935,-0.346) (1.950,-0.314) (1.963,-0.282) (1.975,-0.249) (1.985,-0.215) (1.994,-0.181) (2.001,-0.146) (2.006,-0.112) (2.011,-0.077) (2.013,-0.042) (2.014,-0.007) (2.014,0.028) (2.012,0.063) (2.008,0.098) (2.003,0.133) (1.997,0.168) (1.989,0.202) (1.979,0.236) (1.968,0.269) (1.955,0.302) (1.941,0.334) (1.925,0.365) (1.908,0.396) (1.889,0.426) (1.869,0.454) (1.847,0.482) (1.824,0.508) (1.799,0.533) (1.772,0.556) (1.745,0.577) (1.715,0.597) (1.685,0.614) (1.653,0.629) (1.620,0.641) (1.586,0.650) (1.551,0.656) (1.516,0.658) (1.481,0.657) (1.446,0.652) (1.412,0.643) (1.379,0.632) (1.347,0.617) (1.317,0.600) (1.288,0.580) (1.260,0.558) (1.234,0.534) (1.210,0.509) (1.187,0.482) (1.166,0.454) (1.147,0.425) (1.129,0.395) (1.113,0.363) (1.098,0.331) (1.085,0.299) (1.073,0.266) (1.063,0.232) (1.054,0.198) (1.047,0.164) (1.041,0.129) (1.036,0.094) (1.033,0.059) (1.031,0.024) (1.031,-0.011) (1.032,-0.046) (1.035,-0.081) (1.039,-0.116) (1.044,-0.151) (1.051,-0.185) (1.059,-0.219) (1.069,-0.253) (1.080,-0.286) (1.093,-0.319) (1.107,-0.351) (1.123,-0.383) (1.140,-0.413) (1.159,-0.443) (1.179,-0.472) (1.201,-0.499) (1.225,-0.525) (1.250,-0.549) (1.277,-0.572) (1.305,-0.593) (1.335,-0.611) (1.367,-0.627) (1.399,-0.639) (1.433,-0.649) (1.468,-0.655)};
\addplot[equipot] coordinates {(1.468,-0.532) (1.503,-0.537) (1.537,-0.532) (1.570,-0.519) (1.600,-0.501) (1.629,-0.480) (1.654,-0.456) (1.678,-0.430) (1.700,-0.402) (1.720,-0.373) (1.738,-0.342) (1.754,-0.311) (1.769,-0.279) (1.782,-0.246) (1.793,-0.212) (1.802,-0.178) (1.809,-0.144) (1.815,-0.109) (1.820,-0.074) (1.822,-0.039) (1.823,-0.004) (1.823,0.032) (1.820,0.067) (1.816,0.102) (1.811,0.137) (1.804,0.171) (1.795,0.205) (1.784,0.239) (1.772,0.272) (1.757,0.304) (1.742,0.336) (1.724,0.367) (1.704,0.396) (1.683,0.424) (1.660,0.450) (1.634,0.475) (1.606,0.497) (1.577,0.515) (1.544,0.529) (1.510,0.537) (1.475,0.534) (1.442,0.522) (1.411,0.504) (1.383,0.483) (1.357,0.459) (1.334,0.432) (1.313,0.404) (1.293,0.374) (1.276,0.344) (1.260,0.312) (1.247,0.280) (1.234,0.246) (1.224,0.213) (1.215,0.179) (1.208,0.144) (1.202,0.109) (1.198,0.074) (1.196,0.039) (1.195,0.004) (1.196,-0.032) (1.198,-0.067) (1.201,-0.102) (1.207,-0.137) (1.214,-0.171) (1.222,-0.206) (1.232,-0.240) (1.244,-0.273) (1.258,-0.305) (1.273,-0.337) (1.290,-0.368) (1.309,-0.398) (1.330,-0.426) (1.352,-0.453) (1.378,-0.478) (1.405,-0.500) (1.435,-0.519) (1.468,-0.532)};
\addplot[equipot] coordinates {(1.500,-0.482) (1.526,-0.459) (1.550,-0.434) (1.572,-0.406) (1.592,-0.378) (1.610,-0.347) (1.627,-0.316) (1.641,-0.284) (1.654,-0.252) (1.665,-0.218) (1.674,-0.185) (1.682,-0.150) (1.688,-0.116) (1.692,-0.081) (1.695,-0.046) (1.697,-0.011) (1.696,0.024) (1.694,0.059) (1.691,0.094) (1.686,0.129) (1.679,0.163) (1.671,0.197) (1.661,0.231) (1.649,0.264) (1.636,0.297) (1.621,0.328) (1.604,0.359) (1.585,0.389) (1.564,0.417) (1.541,0.444) (1.516,0.468) (1.490,0.474) (1.465,0.449) (1.442,0.422) (1.422,0.394) (1.403,0.364) (1.386,0.333) (1.372,0.301) (1.358,0.269) (1.347,0.235) (1.337,0.202) (1.329,0.168) (1.323,0.133) (1.317,0.098) (1.314,0.063) (1.312,0.028) (1.312,-0.007) (1.313,-0.042) (1.315,-0.077) (1.319,-0.112) (1.325,-0.146) (1.332,-0.181) (1.341,-0.215) (1.351,-0.248) (1.363,-0.281) (1.377,-0.313) (1.392,-0.345) (1.410,-0.375) (1.429,-0.405) (1.451,-0.432) (1.474,-0.458) (1.500,-0.482)};
\addplot[equipot] coordinates {(1.500,-0.395) (1.516,-0.364) (1.531,-0.332) (1.545,-0.300) (1.557,-0.267) (1.568,-0.234) (1.577,-0.200) (1.585,-0.166) (1.591,-0.131) (1.596,-0.097) (1.599,-0.062) (1.601,-0.027) (1.601,0.008) (1.600,0.043) (1.598,0.078) (1.594,0.113) (1.588,0.147) (1.581,0.182) (1.573,0.216) (1.563,0.249) (1.552,0.282) (1.539,0.315) (1.524,0.347) (1.508,0.378) (1.494,0.380) (1.477,0.349) (1.463,0.317) (1.451,0.284) (1.439,0.251) (1.430,0.217) (1.422,0.183) (1.415,0.149) (1.410,0.114) (1.406,0.079) (1.403,0.044) (1.402,0.009) (1.402,-0.026) (1.404,-0.061) (1.407,-0.095) (1.412,-0.130) (1.418,-0.165) (1.425,-0.199) (1.434,-0.233) (1.444,-0.266) (1.456,-0.299) (1.470,-0.332) (1.484,-0.363) (1.500,-0.395)};
\addplot[lineaE, flechaen=0.375] coordinates {(1.500,0.112) (1.535,0.115) (1.571,0.120) (1.606,0.125) (1.641,0.131) (1.676,0.137) (1.711,0.143) (1.746,0.150) (1.780,0.158) (1.815,0.165) (1.850,0.173) (1.884,0.182) (1.919,0.190) (1.953,0.199) (1.988,0.208) (2.022,0.218) (2.056,0.227) (2.090,0.237) (2.125,0.247) (2.159,0.257) (2.193,0.267) (2.227,0.277) (2.261,0.288) (2.295,0.298) (2.329,0.309) (2.362,0.320) (2.396,0.331) (2.430,0.342) (2.464,0.353) (2.498,0.364) (2.531,0.376) (2.565,0.387) (2.599,0.398) (2.632,0.410) (2.666,0.422) (2.699,0.433) (2.733,0.445) (2.766,0.457) (2.800,0.469) (2.833,0.481) (2.867,0.493) (2.900,0.505) (2.934,0.517) (2.967,0.529) (3.000,0.542) (3.034,0.554) (3.067,0.567) (3.100,0.579) (3.134,0.592) (3.167,0.604) (3.200,0.617)};
\addplot[lineaE, flechaen=0.625] coordinates {(-3.200,0.617) (-3.167,0.604) (-3.134,0.592) (-3.100,0.579) (-3.067,0.567) (-3.034,0.554) (-3.000,0.542) (-2.967,0.529) (-2.934,0.517) (-2.900,0.505) (-2.867,0.493) (-2.833,0.481) (-2.800,0.469) (-2.766,0.457) (-2.733,0.445) (-2.699,0.433) (-2.666,0.422) (-2.632,0.410) (-2.599,0.398) (-2.565,0.387) (-2.531,0.376) (-2.498,0.364) (-2.464,0.353) (-2.430,0.342) (-2.396,0.331) (-2.362,0.320) (-2.329,0.309) (-2.295,0.298) (-2.261,0.288) (-2.227,0.277) (-2.193,0.267) (-2.159,0.257) (-2.125,0.247) (-2.090,0.237) (-2.056,0.227) (-2.022,0.218) (-1.988,0.208) (-1.953,0.199) (-1.919,0.190) (-1.884,0.182) (-1.850,0.173) (-1.815,0.165) (-1.780,0.158) (-1.746,0.150) (-1.711,0.143) (-1.676,0.137) (-1.641,0.131) (-1.606,0.125) (-1.571,0.120) (-1.535,0.115) (-1.500,0.112)};
\addplot[lineaE, flechaen=0.348] coordinates {(1.500,0.218) (1.534,0.225) (1.568,0.234) (1.601,0.245) (1.635,0.256) (1.667,0.269) (1.700,0.282) (1.732,0.296) (1.764,0.311) (1.795,0.327) (1.826,0.343) (1.857,0.360) (1.888,0.377) (1.918,0.394) (1.948,0.412) (1.978,0.431) (2.008,0.449) (2.037,0.468) (2.066,0.487) (2.096,0.507) (2.125,0.526) (2.154,0.546) (2.182,0.566) (2.211,0.587) (2.239,0.607) (2.268,0.628) (2.296,0.648) (2.324,0.669) (2.352,0.690) (2.380,0.711) (2.408,0.733) (2.436,0.754) (2.463,0.776) (2.491,0.797) (2.518,0.819) (2.546,0.841) (2.573,0.863) (2.600,0.885) (2.627,0.908) (2.654,0.930) (2.681,0.953) (2.708,0.975) (2.734,0.998) (2.761,1.021) (2.788,1.044) (2.814,1.067) (2.840,1.090) (2.867,1.113) (2.893,1.136) (2.919,1.160) (2.945,1.183) (2.971,1.207) (2.997,1.231) (3.022,1.254) (3.048,1.278) (3.074,1.302) (3.099,1.326) (3.124,1.351) (3.150,1.375) (3.175,1.399) (3.200,1.424)};
\addplot[lineaE, flechaen=0.652] coordinates {(-3.200,1.424) (-3.175,1.399) (-3.150,1.375) (-3.124,1.351) (-3.099,1.326) (-3.074,1.302) (-3.048,1.278) (-3.022,1.254) (-2.997,1.231) (-2.971,1.207) (-2.945,1.183) (-2.919,1.160) (-2.893,1.136) (-2.867,1.113) (-2.840,1.090) (-2.814,1.067) (-2.788,1.044) (-2.761,1.021) (-2.734,0.998) (-2.708,0.975) (-2.681,0.953) (-2.654,0.930) (-2.627,0.908) (-2.600,0.885) (-2.573,0.863) (-2.546,0.841) (-2.518,0.819) (-2.491,0.797) (-2.463,0.776) (-2.436,0.754) (-2.408,0.733) (-2.380,0.711) (-2.352,0.690) (-2.324,0.669) (-2.296,0.648) (-2.268,0.628) (-2.239,0.607) (-2.211,0.587) (-2.182,0.566) (-2.154,0.546) (-2.125,0.526) (-2.096,0.507) (-2.066,0.487) (-2.037,0.468) (-2.008,0.449) (-1.978,0.431) (-1.948,0.412) (-1.918,0.394) (-1.888,0.377) (-1.857,0.360) (-1.826,0.343) (-1.795,0.327) (-1.764,0.311) (-1.732,0.296) (-1.700,0.282) (-1.667,0.269) (-1.635,0.256) (-1.601,0.245) (-1.568,0.234) (-1.534,0.225) (-1.500,0.218)};
\addplot[lineaE, flechaen=0.508] coordinates {(1.500,0.325) (1.533,0.336) (1.565,0.352) (1.595,0.369) (1.624,0.389) (1.653,0.410) (1.680,0.432) (1.706,0.456) (1.732,0.480) (1.757,0.505) (1.781,0.531) (1.804,0.557) (1.827,0.584) (1.850,0.611) (1.872,0.639) (1.894,0.666) (1.915,0.694) (1.936,0.723) (1.957,0.751) (1.978,0.780) (1.998,0.809) (2.018,0.838) (2.038,0.867) (2.057,0.896) (2.076,0.926) (2.095,0.956) (2.114,0.986) (2.133,1.016) (2.151,1.046) (2.169,1.076) (2.187,1.106) (2.205,1.137) (2.223,1.167) (2.240,1.198) (2.257,1.229) (2.274,1.260) (2.291,1.291) (2.308,1.322) (2.324,1.353) (2.340,1.385) (2.356,1.416) (2.372,1.448) (2.387,1.479) (2.402,1.511) (2.418,1.543) (2.432,1.575) (2.447,1.607) (2.462,1.639) (2.476,1.672) (2.490,1.704) (2.504,1.736) (2.517,1.769) (2.531,1.802) (2.544,1.834) (2.557,1.867) (2.570,1.900)};
\addplot[lineaE, flechaen=0.492] coordinates {(-2.570,1.900) (-2.557,1.867) (-2.544,1.834) (-2.531,1.802) (-2.517,1.769) (-2.504,1.736) (-2.490,1.704) (-2.476,1.672) (-2.462,1.639) (-2.447,1.607) (-2.432,1.575) (-2.418,1.543) (-2.402,1.511) (-2.387,1.479) (-2.372,1.448) (-2.356,1.416) (-2.340,1.385) (-2.324,1.353) (-2.308,1.322) (-2.291,1.291) (-2.274,1.260) (-2.257,1.229) (-2.240,1.198) (-2.223,1.167) (-2.205,1.137) (-2.187,1.106) (-2.169,1.076) (-2.151,1.046) (-2.133,1.016) (-2.114,0.986) (-2.095,0.956) (-2.076,0.926) (-2.057,0.896) (-2.038,0.867) (-2.018,0.838) (-1.998,0.809) (-1.978,0.780) (-1.957,0.751) (-1.936,0.723) (-1.915,0.694) (-1.894,0.666) (-1.872,0.639) (-1.850,0.611) (-1.827,0.584) (-1.804,0.557) (-1.781,0.531) (-1.757,0.505) (-1.732,0.480) (-1.706,0.456) (-1.680,0.432) (-1.653,0.410) (-1.624,0.389) (-1.595,0.369) (-1.565,0.352) (-1.533,0.336) (-1.500,0.325)};
\addplot[lineaE, flechaen=0.500] coordinates {(1.500,0.433) (1.529,0.451) (1.553,0.477) (1.572,0.507) (1.588,0.538) (1.602,0.571) (1.615,0.604) (1.627,0.638) (1.638,0.671) (1.649,0.705) (1.658,0.739) (1.668,0.774) (1.677,0.808) (1.685,0.842) (1.693,0.877) (1.701,0.912) (1.708,0.946) (1.715,0.981) (1.722,1.016) (1.728,1.051) (1.734,1.086) (1.739,1.121) (1.745,1.156) (1.750,1.191) (1.754,1.227) (1.758,1.262) (1.762,1.297) (1.766,1.333) (1.769,1.368) (1.772,1.403) (1.774,1.439) (1.776,1.474) (1.778,1.510) (1.779,1.545) (1.780,1.581) (1.781,1.616) (1.781,1.652) (1.781,1.687) (1.780,1.723) (1.779,1.758) (1.778,1.794) (1.776,1.829) (1.774,1.865) (1.772,1.900)};
\addplot[lineaE, flechaen=0.500] coordinates {(-1.772,1.900) (-1.774,1.865) (-1.776,1.829) (-1.778,1.794) (-1.779,1.758) (-1.780,1.723) (-1.781,1.687) (-1.781,1.652) (-1.781,1.616) (-1.780,1.581) (-1.779,1.545) (-1.778,1.510) (-1.776,1.474) (-1.774,1.439) (-1.772,1.403) (-1.769,1.368) (-1.766,1.333) (-1.762,1.297) (-1.758,1.262) (-1.754,1.227) (-1.750,1.191) (-1.745,1.156) (-1.739,1.121) (-1.734,1.086) (-1.728,1.051) (-1.722,1.016) (-1.715,0.981) (-1.708,0.946) (-1.701,0.912) (-1.693,0.877) (-1.685,0.842) (-1.677,0.808) (-1.668,0.774) (-1.658,0.739) (-1.649,0.705) (-1.638,0.671) (-1.627,0.638) (-1.615,0.604) (-1.602,0.571) (-1.588,0.538) (-1.572,0.507) (-1.553,0.477) (-1.529,0.451) (-1.500,0.433)};
\addplot[lineaE, flechaen=0.500] coordinates {(1.472,0.500) (1.461,0.534) (1.451,0.568) (1.442,0.602) (1.433,0.637) (1.425,0.672) (1.416,0.706) (1.408,0.741) (1.399,0.775) (1.391,0.810) (1.382,0.845) (1.372,0.879) (1.363,0.913) (1.353,0.948) (1.343,0.982) (1.332,1.016) (1.322,1.050) (1.310,1.084) (1.299,1.117) (1.287,1.151) (1.274,1.184) (1.261,1.218) (1.248,1.251) (1.234,1.283) (1.220,1.316) (1.205,1.349) (1.189,1.381) (1.174,1.413) (1.157,1.444) (1.140,1.476) (1.123,1.507) (1.105,1.538) (1.086,1.568) (1.067,1.598) (1.047,1.628) (1.027,1.657) (1.006,1.686) (0.985,1.715) (0.963,1.743) (0.940,1.770) (0.917,1.797) (0.893,1.824) (0.869,1.850) (0.844,1.875) (0.818,1.900)};
\addplot[lineaE, flechaen=0.500] coordinates {(-0.818,1.900) (-0.844,1.875) (-0.869,1.850) (-0.893,1.824) (-0.917,1.797) (-0.940,1.770) (-0.963,1.743) (-0.985,1.715) (-1.006,1.686) (-1.027,1.657) (-1.047,1.628) (-1.067,1.598) (-1.086,1.568) (-1.105,1.538) (-1.123,1.507) (-1.140,1.476) (-1.157,1.444) (-1.174,1.413) (-1.189,1.381) (-1.205,1.349) (-1.220,1.316) (-1.234,1.283) (-1.248,1.251) (-1.261,1.218) (-1.274,1.184) (-1.287,1.151) (-1.299,1.117) (-1.310,1.084) (-1.322,1.050) (-1.332,1.016) (-1.343,0.982) (-1.353,0.948) (-1.363,0.913) (-1.372,0.879) (-1.382,0.845) (-1.391,0.810) (-1.399,0.775) (-1.408,0.741) (-1.416,0.706) (-1.425,0.672) (-1.433,0.637) (-1.442,0.602) (-1.451,0.568) (-1.461,0.534) (-1.472,0.500)};
\addplot[lineaE, flechaen=0.686] coordinates {(1.196,0.500) (1.173,0.525) (1.147,0.549) (1.121,0.572) (1.094,0.596) (1.068,0.619) (1.042,0.642) (1.015,0.666) (0.989,0.689) (0.962,0.712) (0.935,0.734) (0.908,0.757) (0.880,0.779) (0.853,0.800) (0.825,0.822) (0.796,0.843) (0.768,0.863) (0.739,0.883) (0.709,0.903) (0.680,0.922) (0.650,0.940) (0.619,0.958) (0.589,0.975) (0.558,0.991) (0.526,1.007) (0.494,1.022) (0.462,1.036) (0.429,1.050) (0.396,1.062) (0.363,1.073) (0.330,1.084) (0.296,1.094) (0.262,1.102) (0.227,1.110) (0.193,1.117) (0.158,1.122) (0.123,1.127) (0.088,1.130) (0.053,1.132) (0.018,1.133) (-0.018,1.133) (-0.053,1.132) (-0.088,1.130) (-0.123,1.127) (-0.158,1.122) (-0.193,1.117) (-0.227,1.110) (-0.262,1.102) (-0.296,1.094) (-0.330,1.084) (-0.363,1.073) (-0.396,1.062) (-0.429,1.050) (-0.462,1.036) (-0.494,1.022) (-0.526,1.007) (-0.558,0.991) (-0.589,0.975) (-0.619,0.958) (-0.650,0.940) (-0.680,0.922) (-0.709,0.903) (-0.739,0.883) (-0.768,0.863) (-0.796,0.843) (-0.825,0.822) (-0.853,0.800) (-0.880,0.779) (-0.908,0.757) (-0.935,0.734) (-0.962,0.712) (-0.989,0.689) (-1.015,0.666) (-1.042,0.642) (-1.068,0.619) (-1.094,0.596) (-1.121,0.572) (-1.147,0.549) (-1.173,0.525) (-1.196,0.500)};
\addplot[lineaE, flechaen=0.336] coordinates {(1.500,0.131) (1.465,0.135) (1.430,0.141) (1.396,0.147) (1.361,0.153) (1.326,0.160) (1.292,0.167) (1.258,0.175) (1.223,0.183) (1.189,0.191) (1.155,0.200) (1.121,0.208) (1.087,0.217) (1.053,0.227) (1.019,0.236) (0.985,0.245) (0.951,0.255) (0.917,0.264) (0.883,0.273) (0.849,0.283) (0.815,0.292) (0.781,0.301) (0.747,0.310) (0.713,0.319) (0.679,0.327) (0.644,0.335) (0.610,0.344) (0.576,0.351) (0.541,0.359) (0.507,0.366) (0.472,0.373) (0.438,0.379) (0.403,0.385) (0.368,0.391) (0.333,0.396) (0.299,0.401) (0.264,0.405) (0.229,0.409) (0.193,0.412) (0.158,0.414) (0.123,0.417) (0.088,0.418) (0.053,0.419) (0.018,0.420) (-0.018,0.420) (-0.053,0.419) (-0.088,0.418) (-0.123,0.417) (-0.158,0.414) (-0.193,0.412) (-0.229,0.409) (-0.264,0.405) (-0.299,0.401) (-0.333,0.396) (-0.368,0.391) (-0.403,0.385) (-0.438,0.379) (-0.472,0.373) (-0.507,0.366) (-0.541,0.359) (-0.576,0.351) (-0.610,0.344) (-0.644,0.335) (-0.679,0.327) (-0.713,0.319) (-0.747,0.310) (-0.781,0.301) (-0.815,0.292) (-0.849,0.283) (-0.883,0.273) (-0.917,0.264) (-0.951,0.255) (-0.985,0.245) (-1.019,0.236) (-1.053,0.227) (-1.087,0.217) (-1.121,0.208) (-1.155,0.200) (-1.189,0.191) (-1.223,0.183) (-1.258,0.175) (-1.292,0.167) (-1.326,0.160) (-1.361,0.153) (-1.396,0.147) (-1.430,0.141) (-1.465,0.135) (-1.500,0.131)};
\addplot[lineaE, flechaen=0.665] coordinates {(1.500,0.101) (1.465,0.104) (1.430,0.108) (1.395,0.113) (1.360,0.117) (1.325,0.123) (1.290,0.128) (1.255,0.134) (1.220,0.140) (1.186,0.147) (1.151,0.153) (1.116,0.160) (1.082,0.167) (1.047,0.174) (1.012,0.181) (0.978,0.188) (0.943,0.195) (0.909,0.203) (0.874,0.210) (0.840,0.217) (0.805,0.224) (0.770,0.231) (0.736,0.238) (0.701,0.244) (0.666,0.251) (0.632,0.257) (0.597,0.263) (0.562,0.269) (0.527,0.275) (0.492,0.280) (0.457,0.286) (0.422,0.290) (0.387,0.295) (0.352,0.299) (0.317,0.303) (0.282,0.307) (0.247,0.310) (0.212,0.312) (0.176,0.315) (0.141,0.317) (0.106,0.318) (0.071,0.319) (0.035,0.320) (0.000,0.320) (-0.035,0.320) (-0.071,0.319) (-0.106,0.318) (-0.141,0.317) (-0.176,0.315) (-0.212,0.312) (-0.247,0.310) (-0.282,0.307) (-0.317,0.303) (-0.352,0.299) (-0.387,0.295) (-0.422,0.290) (-0.457,0.286) (-0.492,0.280) (-0.527,0.275) (-0.562,0.269) (-0.597,0.263) (-0.632,0.257) (-0.666,0.251) (-0.701,0.244) (-0.736,0.238) (-0.770,0.231) (-0.805,0.224) (-0.840,0.217) (-0.874,0.210) (-0.909,0.203) (-0.943,0.195) (-0.978,0.188) (-1.012,0.181) (-1.047,0.174) (-1.082,0.167) (-1.116,0.160) (-1.151,0.153) (-1.186,0.147) (-1.220,0.140) (-1.255,0.134) (-1.290,0.128) (-1.325,0.123) (-1.360,0.117) (-1.395,0.113) (-1.430,0.108) (-1.465,0.104) (-1.500,0.101)};
\addplot[lineaE, flechaen=0.334] coordinates {(1.500,0.070) (1.465,0.072) (1.430,0.075) (1.395,0.078) (1.360,0.081) (1.325,0.085) (1.290,0.088) (1.255,0.092) (1.221,0.097) (1.186,0.101) (1.151,0.105) (1.116,0.110) (1.082,0.115) (1.047,0.120) (1.012,0.124) (0.977,0.129) (0.943,0.134) (0.908,0.139) (0.873,0.144) (0.838,0.149) (0.804,0.154) (0.769,0.159) (0.734,0.163) (0.699,0.168) (0.664,0.172) (0.630,0.177) (0.595,0.181) (0.560,0.185) (0.525,0.189) (0.490,0.193) (0.455,0.196) (0.420,0.200) (0.385,0.203) (0.350,0.206) (0.315,0.208) (0.280,0.211) (0.245,0.213) (0.210,0.215) (0.175,0.216) (0.140,0.218) (0.105,0.219) (0.070,0.219) (0.035,0.220) (-0.000,0.220) (-0.035,0.220) (-0.070,0.219) (-0.105,0.219) (-0.140,0.218) (-0.175,0.216) (-0.210,0.215) (-0.245,0.213) (-0.280,0.211) (-0.315,0.208) (-0.350,0.206) (-0.385,0.203) (-0.420,0.200) (-0.455,0.196) (-0.490,0.193) (-0.525,0.189) (-0.560,0.185) (-0.595,0.181) (-0.630,0.177) (-0.664,0.172) (-0.699,0.168) (-0.734,0.163) (-0.769,0.159) (-0.804,0.154) (-0.838,0.149) (-0.873,0.144) (-0.908,0.139) (-0.943,0.134) (-0.977,0.129) (-1.012,0.124) (-1.047,0.120) (-1.082,0.115) (-1.116,0.110) (-1.151,0.105) (-1.186,0.101) (-1.221,0.097) (-1.255,0.092) (-1.290,0.088) (-1.325,0.085) (-1.360,0.081) (-1.395,0.078) (-1.430,0.075) (-1.465,0.072) (-1.500,0.070)};
\addplot[lineaE, flechaen=0.666] coordinates {(1.500,0.038) (1.465,0.039) (1.429,0.041) (1.394,0.043) (1.359,0.044) (1.323,0.046) (1.288,0.048) (1.253,0.051) (1.218,0.053) (1.182,0.055) (1.147,0.058) (1.112,0.060) (1.077,0.063) (1.041,0.066) (1.006,0.068) (0.971,0.071) (0.935,0.074) (0.900,0.076) (0.865,0.079) (0.830,0.082) (0.794,0.084) (0.759,0.087) (0.724,0.090) (0.689,0.092) (0.653,0.095) (0.618,0.097) (0.583,0.099) (0.548,0.102) (0.512,0.104) (0.477,0.106) (0.442,0.108) (0.406,0.110) (0.371,0.111) (0.336,0.113) (0.300,0.114) (0.265,0.115) (0.230,0.117) (0.194,0.117) (0.159,0.118) (0.124,0.119) (0.088,0.119) (0.053,0.120) (0.018,0.120) (-0.018,0.120) (-0.053,0.120) (-0.088,0.119) (-0.124,0.119) (-0.159,0.118) (-0.194,0.117) (-0.230,0.117) (-0.265,0.115) (-0.300,0.114) (-0.336,0.113) (-0.371,0.111) (-0.406,0.110) (-0.442,0.108) (-0.477,0.106) (-0.512,0.104) (-0.548,0.102) (-0.583,0.099) (-0.618,0.097) (-0.653,0.095) (-0.689,0.092) (-0.724,0.090) (-0.759,0.087) (-0.794,0.084) (-0.830,0.082) (-0.865,0.079) (-0.900,0.076) (-0.935,0.074) (-0.971,0.071) (-1.006,0.068) (-1.041,0.066) (-1.077,0.063) (-1.112,0.060) (-1.147,0.058) (-1.182,0.055) (-1.218,0.053) (-1.253,0.051) (-1.288,0.048) (-1.323,0.046) (-1.359,0.044) (-1.394,0.043) (-1.429,0.041) (-1.465,0.039) (-1.500,0.038)};
\addplot[lineaE, flechaen=0.375] coordinates {(1.500,-0.112) (1.535,-0.115) (1.571,-0.120) (1.606,-0.125) (1.641,-0.131) (1.676,-0.137) (1.711,-0.143) (1.746,-0.150) (1.780,-0.158) (1.815,-0.165) (1.850,-0.173) (1.884,-0.182) (1.919,-0.190) (1.953,-0.199) (1.988,-0.208) (2.022,-0.218) (2.056,-0.227) (2.090,-0.237) (2.125,-0.247) (2.159,-0.257) (2.193,-0.267) (2.227,-0.277) (2.261,-0.288) (2.295,-0.298) (2.329,-0.309) (2.362,-0.320) (2.396,-0.331) (2.430,-0.342) (2.464,-0.353) (2.498,-0.364) (2.531,-0.376) (2.565,-0.387) (2.599,-0.398) (2.632,-0.410) (2.666,-0.422) (2.699,-0.433) (2.733,-0.445) (2.766,-0.457) (2.800,-0.469) (2.833,-0.481) (2.867,-0.493) (2.900,-0.505) (2.934,-0.517) (2.967,-0.529) (3.000,-0.542) (3.034,-0.554) (3.067,-0.567) (3.100,-0.579) (3.134,-0.592) (3.167,-0.604) (3.200,-0.617)};
\addplot[lineaE, flechaen=0.625] coordinates {(-3.200,-0.617) (-3.167,-0.604) (-3.134,-0.592) (-3.100,-0.579) (-3.067,-0.567) (-3.034,-0.554) (-3.000,-0.542) (-2.967,-0.529) (-2.934,-0.517) (-2.900,-0.505) (-2.867,-0.493) (-2.833,-0.481) (-2.800,-0.469) (-2.766,-0.457) (-2.733,-0.445) (-2.699,-0.433) (-2.666,-0.422) (-2.632,-0.410) (-2.599,-0.398) (-2.565,-0.387) (-2.531,-0.376) (-2.498,-0.364) (-2.464,-0.353) (-2.430,-0.342) (-2.396,-0.331) (-2.362,-0.320) (-2.329,-0.309) (-2.295,-0.298) (-2.261,-0.288) (-2.227,-0.277) (-2.193,-0.267) (-2.159,-0.257) (-2.125,-0.247) (-2.090,-0.237) (-2.056,-0.227) (-2.022,-0.218) (-1.988,-0.208) (-1.953,-0.199) (-1.919,-0.190) (-1.884,-0.182) (-1.850,-0.173) (-1.815,-0.165) (-1.780,-0.158) (-1.746,-0.150) (-1.711,-0.143) (-1.676,-0.137) (-1.641,-0.131) (-1.606,-0.125) (-1.571,-0.120) (-1.535,-0.115) (-1.500,-0.112)};
\addplot[lineaE, flechaen=0.348] coordinates {(1.500,-0.218) (1.534,-0.225) (1.568,-0.234) (1.601,-0.245) (1.635,-0.256) (1.667,-0.269) (1.700,-0.282) (1.732,-0.296) (1.764,-0.311) (1.795,-0.327) (1.826,-0.343) (1.857,-0.360) (1.888,-0.377) (1.918,-0.394) (1.948,-0.412) (1.978,-0.431) (2.008,-0.449) (2.037,-0.468) (2.066,-0.487) (2.096,-0.507) (2.125,-0.526) (2.154,-0.546) (2.182,-0.566) (2.211,-0.587) (2.239,-0.607) (2.268,-0.628) (2.296,-0.648) (2.324,-0.669) (2.352,-0.690) (2.380,-0.711) (2.408,-0.733) (2.436,-0.754) (2.463,-0.776) (2.491,-0.797) (2.518,-0.819) (2.546,-0.841) (2.573,-0.863) (2.600,-0.885) (2.627,-0.908) (2.654,-0.930) (2.681,-0.953) (2.708,-0.975) (2.734,-0.998) (2.761,-1.021) (2.788,-1.044) (2.814,-1.067) (2.840,-1.090) (2.867,-1.113) (2.893,-1.136) (2.919,-1.160) (2.945,-1.183) (2.971,-1.207) (2.997,-1.231) (3.022,-1.254) (3.048,-1.278) (3.074,-1.302) (3.099,-1.326) (3.124,-1.351) (3.150,-1.375) (3.175,-1.399) (3.200,-1.424)};
\addplot[lineaE, flechaen=0.652] coordinates {(-3.200,-1.424) (-3.175,-1.399) (-3.150,-1.375) (-3.124,-1.351) (-3.099,-1.326) (-3.074,-1.302) (-3.048,-1.278) (-3.022,-1.254) (-2.997,-1.231) (-2.971,-1.207) (-2.945,-1.183) (-2.919,-1.160) (-2.893,-1.136) (-2.867,-1.113) (-2.840,-1.090) (-2.814,-1.067) (-2.788,-1.044) (-2.761,-1.021) (-2.734,-0.998) (-2.708,-0.975) (-2.681,-0.953) (-2.654,-0.930) (-2.627,-0.908) (-2.600,-0.885) (-2.573,-0.863) (-2.546,-0.841) (-2.518,-0.819) (-2.491,-0.797) (-2.463,-0.776) (-2.436,-0.754) (-2.408,-0.733) (-2.380,-0.711) (-2.352,-0.690) (-2.324,-0.669) (-2.296,-0.648) (-2.268,-0.628) (-2.239,-0.607) (-2.211,-0.587) (-2.182,-0.566) (-2.154,-0.546) (-2.125,-0.526) (-2.096,-0.507) (-2.066,-0.487) (-2.037,-0.468) (-2.008,-0.449) (-1.978,-0.431) (-1.948,-0.412) (-1.918,-0.394) (-1.888,-0.377) (-1.857,-0.360) (-1.826,-0.343) (-1.795,-0.327) (-1.764,-0.311) (-1.732,-0.296) (-1.700,-0.282) (-1.667,-0.269) (-1.635,-0.256) (-1.601,-0.245) (-1.568,-0.234) (-1.534,-0.225) (-1.500,-0.218)};
\addplot[lineaE, flechaen=0.508] coordinates {(1.500,-0.325) (1.533,-0.336) (1.565,-0.352) (1.595,-0.369) (1.624,-0.389) (1.653,-0.410) (1.680,-0.432) (1.706,-0.456) (1.732,-0.480) (1.757,-0.505) (1.781,-0.531) (1.804,-0.557) (1.827,-0.584) (1.850,-0.611) (1.872,-0.639) (1.894,-0.666) (1.915,-0.694) (1.936,-0.723) (1.957,-0.751) (1.978,-0.780) (1.998,-0.809) (2.018,-0.838) (2.038,-0.867) (2.057,-0.896) (2.076,-0.926) (2.095,-0.956) (2.114,-0.986) (2.133,-1.016) (2.151,-1.046) (2.169,-1.076) (2.187,-1.106) (2.205,-1.137) (2.223,-1.167) (2.240,-1.198) (2.257,-1.229) (2.274,-1.260) (2.291,-1.291) (2.308,-1.322) (2.324,-1.353) (2.340,-1.385) (2.356,-1.416) (2.372,-1.448) (2.387,-1.479) (2.402,-1.511) (2.418,-1.543) (2.432,-1.575) (2.447,-1.607) (2.462,-1.639) (2.476,-1.672) (2.490,-1.704) (2.504,-1.736) (2.517,-1.769) (2.531,-1.802) (2.544,-1.834) (2.557,-1.867) (2.570,-1.900)};
\addplot[lineaE, flechaen=0.492] coordinates {(-2.570,-1.900) (-2.557,-1.867) (-2.544,-1.834) (-2.531,-1.802) (-2.517,-1.769) (-2.504,-1.736) (-2.490,-1.704) (-2.476,-1.672) (-2.462,-1.639) (-2.447,-1.607) (-2.432,-1.575) (-2.418,-1.543) (-2.402,-1.511) (-2.387,-1.479) (-2.372,-1.448) (-2.356,-1.416) (-2.340,-1.385) (-2.324,-1.353) (-2.308,-1.322) (-2.291,-1.291) (-2.274,-1.260) (-2.257,-1.229) (-2.240,-1.198) (-2.223,-1.167) (-2.205,-1.137) (-2.187,-1.106) (-2.169,-1.076) (-2.151,-1.046) (-2.133,-1.016) (-2.114,-0.986) (-2.095,-0.956) (-2.076,-0.926) (-2.057,-0.896) (-2.038,-0.867) (-2.018,-0.838) (-1.998,-0.809) (-1.978,-0.780) (-1.957,-0.751) (-1.936,-0.723) (-1.915,-0.694) (-1.894,-0.666) (-1.872,-0.639) (-1.850,-0.611) (-1.827,-0.584) (-1.804,-0.557) (-1.781,-0.531) (-1.757,-0.505) (-1.732,-0.480) (-1.706,-0.456) (-1.680,-0.432) (-1.653,-0.410) (-1.624,-0.389) (-1.595,-0.369) (-1.565,-0.352) (-1.533,-0.336) (-1.500,-0.325)};
\addplot[lineaE, flechaen=0.500] coordinates {(1.500,-0.433) (1.529,-0.451) (1.553,-0.477) (1.572,-0.507) (1.588,-0.538) (1.602,-0.571) (1.615,-0.604) (1.627,-0.638) (1.638,-0.671) (1.649,-0.705) (1.658,-0.739) (1.668,-0.774) (1.677,-0.808) (1.685,-0.842) (1.693,-0.877) (1.701,-0.912) (1.708,-0.946) (1.715,-0.981) (1.722,-1.016) (1.728,-1.051) (1.734,-1.086) (1.739,-1.121) (1.745,-1.156) (1.750,-1.191) (1.754,-1.227) (1.758,-1.262) (1.762,-1.297) (1.766,-1.333) (1.769,-1.368) (1.772,-1.403) (1.774,-1.439) (1.776,-1.474) (1.778,-1.510) (1.779,-1.545) (1.780,-1.581) (1.781,-1.616) (1.781,-1.652) (1.781,-1.687) (1.780,-1.723) (1.779,-1.758) (1.778,-1.794) (1.776,-1.829) (1.774,-1.865) (1.772,-1.900)};
\addplot[lineaE, flechaen=0.500] coordinates {(-1.772,-1.900) (-1.774,-1.865) (-1.776,-1.829) (-1.778,-1.794) (-1.779,-1.758) (-1.780,-1.723) (-1.781,-1.687) (-1.781,-1.652) (-1.781,-1.616) (-1.780,-1.581) (-1.779,-1.545) (-1.778,-1.510) (-1.776,-1.474) (-1.774,-1.439) (-1.772,-1.403) (-1.769,-1.368) (-1.766,-1.333) (-1.762,-1.297) (-1.758,-1.262) (-1.754,-1.227) (-1.750,-1.191) (-1.745,-1.156) (-1.739,-1.121) (-1.734,-1.086) (-1.728,-1.051) (-1.722,-1.016) (-1.715,-0.981) (-1.708,-0.946) (-1.701,-0.912) (-1.693,-0.877) (-1.685,-0.842) (-1.677,-0.808) (-1.668,-0.774) (-1.658,-0.739) (-1.649,-0.705) (-1.638,-0.671) (-1.627,-0.638) (-1.615,-0.604) (-1.602,-0.571) (-1.588,-0.538) (-1.572,-0.507) (-1.553,-0.477) (-1.529,-0.451) (-1.500,-0.433)};
\addplot[lineaE, flechaen=0.500] coordinates {(1.472,-0.500) (1.461,-0.534) (1.451,-0.568) (1.442,-0.602) (1.433,-0.637) (1.425,-0.672) (1.416,-0.706) (1.408,-0.741) (1.399,-0.775) (1.391,-0.810) (1.382,-0.845) (1.372,-0.879) (1.363,-0.913) (1.353,-0.948) (1.343,-0.982) (1.332,-1.016) (1.322,-1.050) (1.310,-1.084) (1.299,-1.117) (1.287,-1.151) (1.274,-1.184) (1.261,-1.218) (1.248,-1.251) (1.234,-1.283) (1.220,-1.316) (1.205,-1.349) (1.189,-1.381) (1.174,-1.413) (1.157,-1.444) (1.140,-1.476) (1.123,-1.507) (1.105,-1.538) (1.086,-1.568) (1.067,-1.598) (1.047,-1.628) (1.027,-1.657) (1.006,-1.686) (0.985,-1.715) (0.963,-1.743) (0.940,-1.770) (0.917,-1.797) (0.893,-1.824) (0.869,-1.850) (0.844,-1.875) (0.818,-1.900)};
\addplot[lineaE, flechaen=0.500] coordinates {(-0.818,-1.900) (-0.844,-1.875) (-0.869,-1.850) (-0.893,-1.824) (-0.917,-1.797) (-0.940,-1.770) (-0.963,-1.743) (-0.985,-1.715) (-1.006,-1.686) (-1.027,-1.657) (-1.047,-1.628) (-1.067,-1.598) (-1.086,-1.568) (-1.105,-1.538) (-1.123,-1.507) (-1.140,-1.476) (-1.157,-1.444) (-1.174,-1.413) (-1.189,-1.381) (-1.205,-1.349) (-1.220,-1.316) (-1.234,-1.283) (-1.248,-1.251) (-1.261,-1.218) (-1.274,-1.184) (-1.287,-1.151) (-1.299,-1.117) (-1.310,-1.084) (-1.322,-1.050) (-1.332,-1.016) (-1.343,-0.982) (-1.353,-0.948) (-1.363,-0.913) (-1.372,-0.879) (-1.382,-0.845) (-1.391,-0.810) (-1.399,-0.775) (-1.408,-0.741) (-1.416,-0.706) (-1.425,-0.672) (-1.433,-0.637) (-1.442,-0.602) (-1.451,-0.568) (-1.461,-0.534) (-1.472,-0.500)};
\addplot[lineaE, flechaen=0.686] coordinates {(1.196,-0.500) (1.173,-0.525) (1.147,-0.549) (1.121,-0.572) (1.094,-0.596) (1.068,-0.619) (1.042,-0.642) (1.015,-0.666) (0.989,-0.689) (0.962,-0.712) (0.935,-0.734) (0.908,-0.757) (0.880,-0.779) (0.853,-0.800) (0.825,-0.822) (0.796,-0.843) (0.768,-0.863) (0.739,-0.883) (0.709,-0.903) (0.680,-0.922) (0.650,-0.940) (0.619,-0.958) (0.589,-0.975) (0.558,-0.991) (0.526,-1.007) (0.494,-1.022) (0.462,-1.036) (0.429,-1.050) (0.396,-1.062) (0.363,-1.073) (0.330,-1.084) (0.296,-1.094) (0.262,-1.102) (0.227,-1.110) (0.193,-1.117) (0.158,-1.122) (0.123,-1.127) (0.088,-1.130) (0.053,-1.132) (0.018,-1.133) (-0.018,-1.133) (-0.053,-1.132) (-0.088,-1.130) (-0.123,-1.127) (-0.158,-1.122) (-0.193,-1.117) (-0.227,-1.110) (-0.262,-1.102) (-0.296,-1.094) (-0.330,-1.084) (-0.363,-1.073) (-0.396,-1.062) (-0.429,-1.050) (-0.462,-1.036) (-0.494,-1.022) (-0.526,-1.007) (-0.558,-0.991) (-0.589,-0.975) (-0.619,-0.958) (-0.650,-0.940) (-0.680,-0.922) (-0.709,-0.903) (-0.739,-0.883) (-0.768,-0.863) (-0.796,-0.843) (-0.825,-0.822) (-0.853,-0.800) (-0.880,-0.779) (-0.908,-0.757) (-0.935,-0.734) (-0.962,-0.712) (-0.989,-0.689) (-1.015,-0.666) (-1.042,-0.642) (-1.068,-0.619) (-1.094,-0.596) (-1.121,-0.572) (-1.147,-0.549) (-1.173,-0.525) (-1.196,-0.500)};
\addplot[lineaE, flechaen=0.336] coordinates {(1.500,-0.131) (1.465,-0.135) (1.430,-0.141) (1.396,-0.147) (1.361,-0.153) (1.326,-0.160) (1.292,-0.167) (1.258,-0.175) (1.223,-0.183) (1.189,-0.191) (1.155,-0.200) (1.121,-0.208) (1.087,-0.217) (1.053,-0.227) (1.019,-0.236) (0.985,-0.245) (0.951,-0.255) (0.917,-0.264) (0.883,-0.273) (0.849,-0.283) (0.815,-0.292) (0.781,-0.301) (0.747,-0.310) (0.713,-0.319) (0.679,-0.327) (0.644,-0.335) (0.610,-0.344) (0.576,-0.351) (0.541,-0.359) (0.507,-0.366) (0.472,-0.373) (0.438,-0.379) (0.403,-0.385) (0.368,-0.391) (0.333,-0.396) (0.299,-0.401) (0.264,-0.405) (0.229,-0.409) (0.193,-0.412) (0.158,-0.414) (0.123,-0.417) (0.088,-0.418) (0.053,-0.419) (0.018,-0.420) (-0.018,-0.420) (-0.053,-0.419) (-0.088,-0.418) (-0.123,-0.417) (-0.158,-0.414) (-0.193,-0.412) (-0.229,-0.409) (-0.264,-0.405) (-0.299,-0.401) (-0.333,-0.396) (-0.368,-0.391) (-0.403,-0.385) (-0.438,-0.379) (-0.472,-0.373) (-0.507,-0.366) (-0.541,-0.359) (-0.576,-0.351) (-0.610,-0.344) (-0.644,-0.335) (-0.679,-0.327) (-0.713,-0.319) (-0.747,-0.310) (-0.781,-0.301) (-0.815,-0.292) (-0.849,-0.283) (-0.883,-0.273) (-0.917,-0.264) (-0.951,-0.255) (-0.985,-0.245) (-1.019,-0.236) (-1.053,-0.227) (-1.087,-0.217) (-1.121,-0.208) (-1.155,-0.200) (-1.189,-0.191) (-1.223,-0.183) (-1.258,-0.175) (-1.292,-0.167) (-1.326,-0.160) (-1.361,-0.153) (-1.396,-0.147) (-1.430,-0.141) (-1.465,-0.135) (-1.500,-0.131)};
\addplot[lineaE, flechaen=0.665] coordinates {(1.500,-0.101) (1.465,-0.104) (1.430,-0.108) (1.395,-0.113) (1.360,-0.117) (1.325,-0.123) (1.290,-0.128) (1.255,-0.134) (1.220,-0.140) (1.186,-0.147) (1.151,-0.153) (1.116,-0.160) (1.082,-0.167) (1.047,-0.174) (1.012,-0.181) (0.978,-0.188) (0.943,-0.195) (0.909,-0.203) (0.874,-0.210) (0.840,-0.217) (0.805,-0.224) (0.770,-0.231) (0.736,-0.238) (0.701,-0.244) (0.666,-0.251) (0.632,-0.257) (0.597,-0.263) (0.562,-0.269) (0.527,-0.275) (0.492,-0.280) (0.457,-0.286) (0.422,-0.290) (0.387,-0.295) (0.352,-0.299) (0.317,-0.303) (0.282,-0.307) (0.247,-0.310) (0.212,-0.312) (0.176,-0.315) (0.141,-0.317) (0.106,-0.318) (0.071,-0.319) (0.035,-0.320) (0.000,-0.320) (-0.035,-0.320) (-0.071,-0.319) (-0.106,-0.318) (-0.141,-0.317) (-0.176,-0.315) (-0.212,-0.312) (-0.247,-0.310) (-0.282,-0.307) (-0.317,-0.303) (-0.352,-0.299) (-0.387,-0.295) (-0.422,-0.290) (-0.457,-0.286) (-0.492,-0.280) (-0.527,-0.275) (-0.562,-0.269) (-0.597,-0.263) (-0.632,-0.257) (-0.666,-0.251) (-0.701,-0.244) (-0.736,-0.238) (-0.770,-0.231) (-0.805,-0.224) (-0.840,-0.217) (-0.874,-0.210) (-0.909,-0.203) (-0.943,-0.195) (-0.978,-0.188) (-1.012,-0.181) (-1.047,-0.174) (-1.082,-0.167) (-1.116,-0.160) (-1.151,-0.153) (-1.186,-0.147) (-1.220,-0.140) (-1.255,-0.134) (-1.290,-0.128) (-1.325,-0.123) (-1.360,-0.117) (-1.395,-0.113) (-1.430,-0.108) (-1.465,-0.104) (-1.500,-0.101)};
\addplot[lineaE, flechaen=0.334] coordinates {(1.500,-0.070) (1.465,-0.072) (1.430,-0.075) (1.395,-0.078) (1.360,-0.081) (1.325,-0.085) (1.290,-0.088) (1.255,-0.092) (1.221,-0.097) (1.186,-0.101) (1.151,-0.105) (1.116,-0.110) (1.082,-0.115) (1.047,-0.120) (1.012,-0.124) (0.977,-0.129) (0.943,-0.134) (0.908,-0.139) (0.873,-0.144) (0.838,-0.149) (0.804,-0.154) (0.769,-0.159) (0.734,-0.163) (0.699,-0.168) (0.664,-0.172) (0.630,-0.177) (0.595,-0.181) (0.560,-0.185) (0.525,-0.189) (0.490,-0.193) (0.455,-0.196) (0.420,-0.200) (0.385,-0.203) (0.350,-0.206) (0.315,-0.208) (0.280,-0.211) (0.245,-0.213) (0.210,-0.215) (0.175,-0.216) (0.140,-0.218) (0.105,-0.219) (0.070,-0.219) (0.035,-0.220) (-0.000,-0.220) (-0.035,-0.220) (-0.070,-0.219) (-0.105,-0.219) (-0.140,-0.218) (-0.175,-0.216) (-0.210,-0.215) (-0.245,-0.213) (-0.280,-0.211) (-0.315,-0.208) (-0.350,-0.206) (-0.385,-0.203) (-0.420,-0.200) (-0.455,-0.196) (-0.490,-0.193) (-0.525,-0.189) (-0.560,-0.185) (-0.595,-0.181) (-0.630,-0.177) (-0.664,-0.172) (-0.699,-0.168) (-0.734,-0.163) (-0.769,-0.159) (-0.804,-0.154) (-0.838,-0.149) (-0.873,-0.144) (-0.908,-0.139) (-0.943,-0.134) (-0.977,-0.129) (-1.012,-0.124) (-1.047,-0.120) (-1.082,-0.115) (-1.116,-0.110) (-1.151,-0.105) (-1.186,-0.101) (-1.221,-0.097) (-1.255,-0.092) (-1.290,-0.088) (-1.325,-0.085) (-1.360,-0.081) (-1.395,-0.078) (-1.430,-0.075) (-1.465,-0.072) (-1.500,-0.070)};
\addplot[lineaE, flechaen=0.666] coordinates {(1.500,-0.038) (1.465,-0.039) (1.429,-0.041) (1.394,-0.043) (1.359,-0.044) (1.323,-0.046) (1.288,-0.048) (1.253,-0.051) (1.218,-0.053) (1.182,-0.055) (1.147,-0.058) (1.112,-0.060) (1.077,-0.063) (1.041,-0.066) (1.006,-0.068) (0.971,-0.071) (0.935,-0.074) (0.900,-0.076) (0.865,-0.079) (0.830,-0.082) (0.794,-0.084) (0.759,-0.087) (0.724,-0.090) (0.689,-0.092) (0.653,-0.095) (0.618,-0.097) (0.583,-0.099) (0.548,-0.102) (0.512,-0.104) (0.477,-0.106) (0.442,-0.108) (0.406,-0.110) (0.371,-0.111) (0.336,-0.113) (0.300,-0.114) (0.265,-0.115) (0.230,-0.117) (0.194,-0.117) (0.159,-0.118) (0.124,-0.119) (0.088,-0.119) (0.053,-0.120) (0.018,-0.120) (-0.018,-0.120) (-0.053,-0.120) (-0.088,-0.119) (-0.124,-0.119) (-0.159,-0.118) (-0.194,-0.117) (-0.230,-0.117) (-0.265,-0.115) (-0.300,-0.114) (-0.336,-0.113) (-0.371,-0.111) (-0.406,-0.110) (-0.442,-0.108) (-0.477,-0.106) (-0.512,-0.104) (-0.548,-0.102) (-0.583,-0.099) (-0.618,-0.097) (-0.653,-0.095) (-0.689,-0.092) (-0.724,-0.090) (-0.759,-0.087) (-0.794,-0.084) (-0.830,-0.082) (-0.865,-0.079) (-0.900,-0.076) (-0.935,-0.074) (-0.971,-0.071) (-1.006,-0.068) (-1.041,-0.066) (-1.077,-0.063) (-1.112,-0.060) (-1.147,-0.058) (-1.182,-0.055) (-1.218,-0.053) (-1.253,-0.051) (-1.288,-0.048) (-1.323,-0.046) (-1.359,-0.044) (-1.394,-0.043) (-1.429,-0.041) (-1.465,-0.039) (-1.500,-0.038)};
\addplot[lineaE, flechaen=0.382] coordinates {(1.500,0.000) (3.200,0.000)};
\addplot[lineaE, flechaen=0.618] coordinates {(-3.200,0.000) (-1.500,0.000)};
\addplot[lineaE, flechaen=0.667] coordinates {(1.500,0.000) (-1.500,0.000)};

  \draw[texto, line width=0.6pt] (axis cs:-1.5,-0.5) rectangle (axis cs:1.5,0.5);
\end{axis}
\end{tikzpicture}
\caption{$\vE(\vr)$}
\end{subfigure}\hfill
\begin{subfigure}[b]{0.325\textwidth}\centering
\begin{tikzpicture}
\begin{axis}[cil]
  \fill[dielec] (axis cs:-1.5,-0.5) rectangle (axis cs:1.5,0.5);
  %% Generado por scripts/lineas_campo.py -- no editar a mano
\addplot[equipot] coordinates {(-1.500,-0.395) (-1.516,-0.364) (-1.531,-0.332) (-1.545,-0.300) (-1.557,-0.267) (-1.568,-0.234) (-1.577,-0.200) (-1.585,-0.166) (-1.591,-0.131) (-1.596,-0.097) (-1.599,-0.062) (-1.601,-0.027) (-1.601,0.008) (-1.600,0.043) (-1.598,0.078) (-1.594,0.113) (-1.588,0.147) (-1.581,0.182) (-1.573,0.216) (-1.563,0.249) (-1.552,0.282) (-1.539,0.315) (-1.524,0.347) (-1.508,0.378) (-1.494,0.380) (-1.477,0.349) (-1.463,0.317) (-1.451,0.284) (-1.439,0.251) (-1.430,0.217) (-1.422,0.183) (-1.415,0.149) (-1.410,0.114) (-1.406,0.079) (-1.403,0.044) (-1.402,0.009) (-1.402,-0.026) (-1.404,-0.061) (-1.407,-0.095) (-1.412,-0.130) (-1.418,-0.165) (-1.425,-0.199) (-1.434,-0.233) (-1.444,-0.266) (-1.456,-0.299) (-1.470,-0.332) (-1.484,-0.363) (-1.500,-0.395)};
\addplot[equipot] coordinates {(-1.500,-0.482) (-1.526,-0.459) (-1.550,-0.434) (-1.572,-0.406) (-1.592,-0.378) (-1.610,-0.347) (-1.627,-0.316) (-1.641,-0.284) (-1.654,-0.252) (-1.665,-0.218) (-1.674,-0.185) (-1.682,-0.150) (-1.688,-0.116) (-1.692,-0.081) (-1.695,-0.046) (-1.697,-0.011) (-1.696,0.024) (-1.694,0.059) (-1.691,0.094) (-1.686,0.129) (-1.679,0.163) (-1.671,0.197) (-1.661,0.231) (-1.649,0.264) (-1.636,0.297) (-1.621,0.328) (-1.604,0.359) (-1.585,0.389) (-1.564,0.417) (-1.541,0.444) (-1.516,0.468) (-1.490,0.474) (-1.465,0.449) (-1.442,0.422) (-1.422,0.394) (-1.403,0.364) (-1.386,0.333) (-1.372,0.301) (-1.358,0.269) (-1.347,0.235) (-1.337,0.202) (-1.329,0.168) (-1.323,0.133) (-1.317,0.098) (-1.314,0.063) (-1.312,0.028) (-1.312,-0.007) (-1.313,-0.042) (-1.315,-0.077) (-1.319,-0.112) (-1.325,-0.146) (-1.332,-0.181) (-1.341,-0.215) (-1.351,-0.248) (-1.363,-0.281) (-1.377,-0.313) (-1.392,-0.345) (-1.410,-0.375) (-1.429,-0.405) (-1.451,-0.432) (-1.474,-0.458) (-1.500,-0.482)};
\addplot[equipot] coordinates {(-1.548,-0.528) (-1.580,-0.514) (-1.610,-0.495) (-1.637,-0.472) (-1.662,-0.448) (-1.685,-0.421) (-1.707,-0.393) (-1.726,-0.363) (-1.743,-0.333) (-1.759,-0.301) (-1.773,-0.269) (-1.785,-0.236) (-1.796,-0.202) (-1.804,-0.168) (-1.812,-0.133) (-1.817,-0.098) (-1.821,-0.063) (-1.823,-0.028) (-1.823,0.007) (-1.822,0.043) (-1.819,0.078) (-1.815,0.113) (-1.809,0.148) (-1.801,0.182) (-1.791,0.216) (-1.780,0.250) (-1.767,0.282) (-1.753,0.315) (-1.736,0.346) (-1.718,0.376) (-1.698,0.405) (-1.676,0.432) (-1.652,0.458) (-1.626,0.482) (-1.597,0.503) (-1.567,0.520) (-1.534,0.533) (-1.499,0.537) (-1.464,0.531) (-1.432,0.517) (-1.402,0.498) (-1.375,0.475) (-1.350,0.450) (-1.327,0.423) (-1.307,0.395) (-1.288,0.365) (-1.271,0.334) (-1.256,0.302) (-1.243,0.269) (-1.231,0.236) (-1.221,0.202) (-1.213,0.168) (-1.206,0.133) (-1.201,0.098) (-1.197,0.063) (-1.196,0.028) (-1.195,-0.007) (-1.196,-0.043) (-1.199,-0.078) (-1.203,-0.113) (-1.209,-0.148) (-1.216,-0.182) (-1.225,-0.216) (-1.236,-0.250) (-1.248,-0.283) (-1.262,-0.316) (-1.278,-0.347) (-1.295,-0.378) (-1.315,-0.407) (-1.337,-0.435) (-1.360,-0.461) (-1.386,-0.485) (-1.414,-0.506) (-1.445,-0.523) (-1.478,-0.535) (-1.514,-0.536) (-1.548,-0.528)};
\addplot[equipot] coordinates {(-1.548,-0.656) (-1.583,-0.651) (-1.617,-0.642) (-1.650,-0.630) (-1.682,-0.616) (-1.712,-0.599) (-1.742,-0.579) (-1.770,-0.558) (-1.796,-0.535) (-1.821,-0.510) (-1.845,-0.484) (-1.867,-0.457) (-1.887,-0.428) (-1.906,-0.399) (-1.924,-0.368) (-1.940,-0.337) (-1.954,-0.305) (-1.967,-0.272) (-1.978,-0.239) (-1.988,-0.205) (-1.996,-0.171) (-2.003,-0.137) (-2.008,-0.102) (-2.011,-0.067) (-2.014,-0.032) (-2.014,0.003) (-2.013,0.039) (-2.011,0.074) (-2.007,0.108) (-2.001,0.143) (-1.994,0.178) (-1.986,0.212) (-1.976,0.245) (-1.964,0.279) (-1.951,0.311) (-1.937,0.343) (-1.920,0.374) (-1.903,0.405) (-1.883,0.434) (-1.863,0.462) (-1.840,0.490) (-1.817,0.515) (-1.791,0.540) (-1.764,0.562) (-1.736,0.583) (-1.707,0.602) (-1.676,0.619) (-1.643,0.633) (-1.610,0.644) (-1.576,0.652) (-1.541,0.657) (-1.506,0.658) (-1.471,0.656) (-1.436,0.650) (-1.403,0.640) (-1.370,0.628) (-1.338,0.612) (-1.308,0.594) (-1.279,0.574) (-1.252,0.552) (-1.227,0.527) (-1.203,0.501) (-1.181,0.474) (-1.160,0.446) (-1.141,0.416) (-1.124,0.386) (-1.108,0.354) (-1.094,0.322) (-1.081,0.289) (-1.070,0.256) (-1.060,0.222) (-1.052,0.188) (-1.045,0.154) (-1.039,0.119) (-1.035,0.084) (-1.032,0.049) (-1.031,0.014) (-1.031,-0.021) (-1.033,-0.056) (-1.036,-0.091) (-1.040,-0.126) (-1.046,-0.161) (-1.053,-0.195) (-1.062,-0.229) (-1.072,-0.263) (-1.084,-0.296) (-1.097,-0.328) (-1.111,-0.360) (-1.127,-0.392) (-1.145,-0.422) (-1.164,-0.451) (-1.185,-0.480) (-1.208,-0.507) (-1.232,-0.532) (-1.257,-0.556) (-1.285,-0.578) (-1.314,-0.598) (-1.344,-0.616) (-1.376,-0.631) (-1.409,-0.642) (-1.443,-0.651) (-1.478,-0.656) (-1.513,-0.658) (-1.548,-0.656)};
\addplot[equipot] coordinates {(-1.660,-0.922) (-1.695,-0.915) (-1.729,-0.907) (-1.763,-0.898) (-1.796,-0.887) (-1.829,-0.874) (-1.861,-0.860) (-1.893,-0.844) (-1.924,-0.827) (-1.954,-0.809) (-1.983,-0.789) (-2.011,-0.769) (-2.039,-0.747) (-2.065,-0.724) (-2.091,-0.699) (-2.115,-0.674) (-2.139,-0.648) (-2.162,-0.621) (-2.183,-0.593) (-2.204,-0.565) (-2.223,-0.535) (-2.241,-0.505) (-2.258,-0.474) (-2.274,-0.443) (-2.289,-0.411) (-2.303,-0.379) (-2.315,-0.346) (-2.327,-0.313) (-2.337,-0.279) (-2.346,-0.245) (-2.354,-0.211) (-2.360,-0.176) (-2.366,-0.141) (-2.370,-0.106) (-2.373,-0.071) (-2.375,-0.036) (-2.376,-0.001) (-2.375,0.034) (-2.373,0.069) (-2.370,0.104) (-2.366,0.139) (-2.361,0.174) (-2.354,0.209) (-2.346,0.243) (-2.337,0.277) (-2.327,0.311) (-2.316,0.344) (-2.304,0.377) (-2.290,0.409) (-2.275,0.441) (-2.259,0.473) (-2.242,0.503) (-2.224,0.534) (-2.205,0.563) (-2.184,0.592) (-2.163,0.620) (-2.140,0.647) (-2.117,0.673) (-2.092,0.698) (-2.067,0.722) (-2.040,0.745) (-2.013,0.767) (-1.985,0.788) (-1.955,0.808) (-1.925,0.826) (-1.895,0.843) (-1.863,0.859) (-1.831,0.873) (-1.798,0.886) (-1.765,0.897) (-1.731,0.907) (-1.697,0.915) (-1.662,0.921) (-1.627,0.926) (-1.592,0.929) (-1.557,0.930) (-1.522,0.930) (-1.487,0.927) (-1.452,0.923) (-1.417,0.917) (-1.383,0.909) (-1.349,0.900) (-1.316,0.889) (-1.283,0.876) (-1.251,0.861) (-1.219,0.846) (-1.189,0.828) (-1.159,0.809) (-1.130,0.789) (-1.103,0.767) (-1.076,0.744) (-1.050,0.720) (-1.026,0.695) (-1.002,0.669) (-0.980,0.642) (-0.959,0.613) (-0.939,0.584) (-0.920,0.555) (-0.902,0.524) (-0.886,0.493) (-0.870,0.462) (-0.856,0.430) (-0.843,0.397) (-0.831,0.364) (-0.820,0.330) (-0.810,0.297) (-0.801,0.263) (-0.794,0.228) (-0.787,0.194) (-0.782,0.159) (-0.778,0.124) (-0.774,0.089) (-0.772,0.054) (-0.771,0.019) (-0.771,-0.017) (-0.772,-0.052) (-0.774,-0.087) (-0.777,-0.122) (-0.782,-0.157) (-0.787,-0.192) (-0.793,-0.226) (-0.801,-0.261) (-0.809,-0.295) (-0.819,-0.329) (-0.830,-0.362) (-0.842,-0.395) (-0.855,-0.428) (-0.869,-0.460) (-0.885,-0.492) (-0.901,-0.523) (-0.919,-0.553) (-0.938,-0.583) (-0.958,-0.612) (-0.979,-0.640) (-1.001,-0.667) (-1.024,-0.693) (-1.049,-0.719) (-1.075,-0.743) (-1.101,-0.766) (-1.129,-0.788) (-1.158,-0.808) (-1.187,-0.827) (-1.218,-0.845) (-1.249,-0.861) (-1.281,-0.875) (-1.314,-0.888) (-1.347,-0.899) (-1.381,-0.909) (-1.415,-0.917) (-1.450,-0.923) (-1.485,-0.927) (-1.520,-0.929) (-1.555,-0.930) (-1.590,-0.929) (-1.625,-0.926) (-1.660,-0.922)};
\addplot[equipot] coordinates {(-3.200,1.227) (-3.178,1.254) (-3.155,1.281) (-3.131,1.307) (-3.107,1.333) (-3.083,1.358) (-3.058,1.382) (-3.033,1.407) (-3.007,1.430) (-2.981,1.454) (-2.954,1.476) (-2.927,1.499) (-2.899,1.520) (-2.871,1.541) (-2.843,1.562) (-2.814,1.582) (-2.785,1.601) (-2.755,1.620) (-2.725,1.639) (-2.695,1.656) (-2.665,1.673) (-2.634,1.690) (-2.602,1.706) (-2.571,1.721) (-2.539,1.735) (-2.507,1.749) (-2.474,1.762) (-2.441,1.775) (-2.408,1.787) (-2.375,1.798) (-2.342,1.808) (-2.308,1.818) (-2.274,1.827) (-2.240,1.835) (-2.206,1.843) (-2.171,1.850) (-2.137,1.856) (-2.102,1.861) (-2.067,1.866) (-2.033,1.869) (-1.998,1.872) (-1.963,1.875) (-1.928,1.876) (-1.893,1.877) (-1.857,1.877) (-1.822,1.876) (-1.787,1.874) (-1.752,1.872) (-1.718,1.868) (-1.683,1.864) (-1.648,1.859) (-1.613,1.854) (-1.579,1.847) (-1.545,1.840) (-1.511,1.832) (-1.477,1.823) (-1.443,1.813) (-1.409,1.803) (-1.376,1.791) (-1.343,1.779) (-1.311,1.766) (-1.278,1.753) (-1.247,1.738) (-1.215,1.723) (-1.184,1.707) (-1.153,1.690) (-1.122,1.673) (-1.093,1.655) (-1.063,1.636) (-1.034,1.616) (-1.005,1.596) (-0.977,1.575) (-0.950,1.553) (-0.923,1.531) (-0.896,1.508) (-0.870,1.484) (-0.845,1.460) (-0.820,1.435) (-0.796,1.410) (-0.772,1.384) (-0.749,1.358) (-0.727,1.331) (-0.705,1.303) (-0.684,1.275) (-0.663,1.247) (-0.643,1.218) (-0.624,1.189) (-0.605,1.159) (-0.588,1.129) (-0.570,1.098) (-0.553,1.067) (-0.537,1.036) (-0.522,1.005) (-0.507,0.973) (-0.493,0.941) (-0.479,0.909) (-0.466,0.876) (-0.454,0.843) (-0.442,0.810) (-0.431,0.777) (-0.420,0.744) (-0.410,0.710) (-0.401,0.676) (-0.392,0.642) (-0.383,0.608) (-0.375,0.574) (-0.368,0.540) (-0.361,0.506) (-0.355,0.471) (-0.349,0.436) (-0.344,0.402) (-0.339,0.367) (-0.334,0.332) (-0.330,0.297) (-0.327,0.263) (-0.324,0.228) (-0.321,0.193) (-0.319,0.158) (-0.317,0.123) (-0.316,0.088) (-0.315,0.053) (-0.315,0.018) (-0.315,-0.018) (-0.315,-0.053) (-0.316,-0.088) (-0.317,-0.123) (-0.319,-0.158) (-0.321,-0.193) (-0.324,-0.228) (-0.327,-0.263) (-0.330,-0.297) (-0.334,-0.332) (-0.339,-0.367) (-0.344,-0.402) (-0.349,-0.436) (-0.355,-0.471) (-0.361,-0.506) (-0.368,-0.540) (-0.375,-0.574) (-0.383,-0.608) (-0.392,-0.642) (-0.401,-0.676) (-0.410,-0.710) (-0.420,-0.744) (-0.431,-0.777) (-0.442,-0.810) (-0.454,-0.843) (-0.466,-0.876) (-0.479,-0.909) (-0.493,-0.941) (-0.507,-0.973) (-0.522,-1.005) (-0.537,-1.036) (-0.553,-1.067) (-0.570,-1.098) (-0.588,-1.129) (-0.605,-1.159) (-0.624,-1.189) (-0.643,-1.218) (-0.663,-1.247) (-0.684,-1.275) (-0.705,-1.303) (-0.727,-1.331) (-0.749,-1.358) (-0.772,-1.384) (-0.796,-1.410) (-0.820,-1.435) (-0.845,-1.460) (-0.870,-1.484) (-0.896,-1.508) (-0.923,-1.531) (-0.950,-1.553) (-0.977,-1.575) (-1.005,-1.596) (-1.034,-1.616) (-1.063,-1.636) (-1.093,-1.655) (-1.122,-1.673) (-1.153,-1.690) (-1.184,-1.707) (-1.215,-1.723) (-1.247,-1.738) (-1.278,-1.753) (-1.311,-1.766) (-1.343,-1.779) (-1.376,-1.791) (-1.409,-1.803) (-1.443,-1.813) (-1.477,-1.823) (-1.511,-1.832) (-1.545,-1.840) (-1.579,-1.847) (-1.613,-1.854) (-1.648,-1.859) (-1.683,-1.864) (-1.718,-1.868) (-1.752,-1.872) (-1.787,-1.874) (-1.822,-1.876) (-1.857,-1.877) (-1.893,-1.877) (-1.928,-1.876) (-1.963,-1.875) (-1.998,-1.872) (-2.033,-1.869) (-2.067,-1.866) (-2.102,-1.861) (-2.137,-1.856) (-2.171,-1.850) (-2.206,-1.843) (-2.240,-1.835) (-2.274,-1.827) (-2.308,-1.818) (-2.342,-1.808) (-2.375,-1.798) (-2.408,-1.787) (-2.441,-1.775) (-2.474,-1.762) (-2.507,-1.749) (-2.539,-1.735) (-2.571,-1.721) (-2.602,-1.706) (-2.634,-1.690) (-2.665,-1.673) (-2.695,-1.656) (-2.725,-1.639) (-2.755,-1.620) (-2.785,-1.601) (-2.814,-1.582) (-2.843,-1.562) (-2.871,-1.541) (-2.899,-1.520) (-2.927,-1.499) (-2.954,-1.476) (-2.981,-1.454) (-3.007,-1.430) (-3.033,-1.407) (-3.058,-1.382) (-3.083,-1.358) (-3.107,-1.333) (-3.131,-1.307) (-3.155,-1.281) (-3.178,-1.254) (-3.200,-1.227)};
\addplot[equipot] coordinates {(3.200,1.227) (3.178,1.254) (3.155,1.281) (3.131,1.307) (3.107,1.333) (3.083,1.358) (3.058,1.382) (3.033,1.407) (3.007,1.430) (2.981,1.454) (2.954,1.476) (2.927,1.499) (2.899,1.520) (2.871,1.541) (2.843,1.562) (2.814,1.582) (2.785,1.601) (2.755,1.620) (2.725,1.639) (2.695,1.656) (2.665,1.673) (2.634,1.690) (2.602,1.706) (2.571,1.721) (2.539,1.735) (2.507,1.749) (2.474,1.762) (2.441,1.775) (2.408,1.787) (2.375,1.798) (2.342,1.808) (2.308,1.818) (2.274,1.827) (2.240,1.835) (2.206,1.843) (2.171,1.850) (2.137,1.856) (2.102,1.861) (2.067,1.866) (2.033,1.869) (1.998,1.872) (1.963,1.875) (1.928,1.876) (1.893,1.877) (1.857,1.877) (1.822,1.876) (1.787,1.874) (1.752,1.872) (1.718,1.868) (1.683,1.864) (1.648,1.859) (1.613,1.854) (1.579,1.847) (1.545,1.840) (1.511,1.832) (1.477,1.823) (1.443,1.813) (1.409,1.803) (1.376,1.791) (1.343,1.779) (1.311,1.766) (1.278,1.753) (1.247,1.738) (1.215,1.723) (1.184,1.707) (1.153,1.690) (1.122,1.673) (1.093,1.655) (1.063,1.636) (1.034,1.616) (1.005,1.596) (0.977,1.575) (0.950,1.553) (0.923,1.531) (0.896,1.508) (0.870,1.484) (0.845,1.460) (0.820,1.435) (0.796,1.410) (0.772,1.384) (0.749,1.358) (0.727,1.331) (0.705,1.303) (0.684,1.275) (0.663,1.247) (0.643,1.218) (0.624,1.189) (0.605,1.159) (0.588,1.129) (0.570,1.098) (0.553,1.067) (0.537,1.036) (0.522,1.005) (0.507,0.973) (0.493,0.941) (0.479,0.909) (0.466,0.876) (0.454,0.843) (0.442,0.810) (0.431,0.777) (0.420,0.744) (0.410,0.710) (0.401,0.676) (0.392,0.642) (0.383,0.608) (0.375,0.574) (0.368,0.540) (0.361,0.506) (0.355,0.471) (0.349,0.436) (0.344,0.402) (0.339,0.367) (0.334,0.332) (0.330,0.297) (0.327,0.263) (0.324,0.228) (0.321,0.193) (0.319,0.158) (0.317,0.123) (0.316,0.088) (0.315,0.053) (0.315,0.018) (0.315,-0.018) (0.315,-0.053) (0.316,-0.088) (0.317,-0.123) (0.319,-0.158) (0.321,-0.193) (0.324,-0.228) (0.327,-0.263) (0.330,-0.297) (0.334,-0.332) (0.339,-0.367) (0.344,-0.402) (0.349,-0.436) (0.355,-0.471) (0.361,-0.506) (0.368,-0.540) (0.375,-0.574) (0.383,-0.608) (0.392,-0.642) (0.401,-0.676) (0.410,-0.710) (0.420,-0.744) (0.431,-0.777) (0.442,-0.810) (0.454,-0.843) (0.466,-0.876) (0.479,-0.909) (0.493,-0.941) (0.507,-0.973) (0.522,-1.005) (0.537,-1.036) (0.553,-1.067) (0.570,-1.098) (0.588,-1.129) (0.605,-1.159) (0.624,-1.189) (0.643,-1.218) (0.663,-1.247) (0.684,-1.275) (0.705,-1.303) (0.727,-1.331) (0.749,-1.358) (0.772,-1.384) (0.796,-1.410) (0.820,-1.435) (0.845,-1.460) (0.870,-1.484) (0.896,-1.508) (0.923,-1.531) (0.950,-1.553) (0.977,-1.575) (1.005,-1.596) (1.034,-1.616) (1.063,-1.636) (1.093,-1.655) (1.122,-1.673) (1.153,-1.690) (1.184,-1.707) (1.215,-1.723) (1.247,-1.738) (1.278,-1.753) (1.311,-1.766) (1.343,-1.779) (1.376,-1.791) (1.409,-1.803) (1.443,-1.813) (1.477,-1.823) (1.511,-1.832) (1.545,-1.840) (1.579,-1.847) (1.613,-1.854) (1.648,-1.859) (1.683,-1.864) (1.718,-1.868) (1.752,-1.872) (1.787,-1.874) (1.822,-1.876) (1.857,-1.877) (1.893,-1.877) (1.928,-1.876) (1.963,-1.875) (1.998,-1.872) (2.033,-1.869) (2.067,-1.866) (2.102,-1.861) (2.137,-1.856) (2.171,-1.850) (2.206,-1.843) (2.240,-1.835) (2.274,-1.827) (2.308,-1.818) (2.342,-1.808) (2.375,-1.798) (2.408,-1.787) (2.441,-1.775) (2.474,-1.762) (2.507,-1.749) (2.539,-1.735) (2.571,-1.721) (2.602,-1.706) (2.634,-1.690) (2.665,-1.673) (2.695,-1.656) (2.725,-1.639) (2.755,-1.620) (2.785,-1.601) (2.814,-1.582) (2.843,-1.562) (2.871,-1.541) (2.899,-1.520) (2.927,-1.499) (2.954,-1.476) (2.981,-1.454) (3.007,-1.430) (3.033,-1.407) (3.058,-1.382) (3.083,-1.358) (3.107,-1.333) (3.131,-1.307) (3.155,-1.281) (3.178,-1.254) (3.200,-1.227)};
\addplot[equipot] coordinates {(1.452,-0.923) (1.487,-0.927) (1.522,-0.930) (1.557,-0.930) (1.592,-0.929) (1.627,-0.926) (1.662,-0.921) (1.697,-0.915) (1.731,-0.907) (1.765,-0.897) (1.798,-0.886) (1.831,-0.873) (1.863,-0.859) (1.894,-0.844) (1.925,-0.826) (1.955,-0.808) (1.984,-0.788) (2.013,-0.768) (2.040,-0.745) (2.067,-0.722) (2.092,-0.698) (2.117,-0.673) (2.140,-0.647) (2.163,-0.620) (2.184,-0.592) (2.205,-0.563) (2.224,-0.534) (2.242,-0.504) (2.259,-0.473) (2.275,-0.441) (2.290,-0.410) (2.304,-0.377) (2.316,-0.344) (2.327,-0.311) (2.337,-0.277) (2.346,-0.243) (2.354,-0.209) (2.361,-0.174) (2.366,-0.139) (2.370,-0.105) (2.373,-0.070) (2.375,-0.034) (2.376,0.001) (2.375,0.036) (2.373,0.071) (2.370,0.106) (2.366,0.141) (2.360,0.176) (2.354,0.210) (2.346,0.245) (2.337,0.279) (2.327,0.312) (2.315,0.346) (2.303,0.379) (2.289,0.411) (2.274,0.443) (2.258,0.474) (2.241,0.505) (2.223,0.535) (2.204,0.565) (2.183,0.593) (2.162,0.621) (2.139,0.648) (2.116,0.674) (2.091,0.699) (2.065,0.723) (2.039,0.747) (2.011,0.769) (1.983,0.789) (1.954,0.809) (1.924,0.827) (1.893,0.844) (1.861,0.860) (1.829,0.874) (1.796,0.887) (1.763,0.898) (1.729,0.907) (1.695,0.915) (1.660,0.922) (1.625,0.926) (1.590,0.929) (1.555,0.930) (1.520,0.929) (1.485,0.927) (1.450,0.923) (1.415,0.917) (1.381,0.909) (1.347,0.899) (1.314,0.888) (1.281,0.875) (1.249,0.861) (1.218,0.845) (1.187,0.827) (1.158,0.808) (1.129,0.788) (1.101,0.766) (1.075,0.743) (1.049,0.719) (1.025,0.694) (1.001,0.667) (0.979,0.640) (0.958,0.612) (0.938,0.583) (0.919,0.553) (0.901,0.523) (0.885,0.492) (0.869,0.460) (0.855,0.428) (0.842,0.395) (0.830,0.362) (0.819,0.329) (0.809,0.295) (0.801,0.261) (0.793,0.226) (0.787,0.192) (0.782,0.157) (0.777,0.122) (0.774,0.087) (0.772,0.052) (0.771,0.017) (0.771,-0.018) (0.772,-0.054) (0.774,-0.089) (0.778,-0.124) (0.782,-0.159) (0.787,-0.193) (0.794,-0.228) (0.801,-0.262) (0.810,-0.296) (0.820,-0.330) (0.831,-0.364) (0.843,-0.397) (0.856,-0.429) (0.870,-0.462) (0.885,-0.493) (0.902,-0.524) (0.920,-0.555) (0.939,-0.584) (0.959,-0.613) (0.980,-0.641) (1.002,-0.669) (1.026,-0.695) (1.050,-0.720) (1.076,-0.744) (1.103,-0.767) (1.130,-0.789) (1.159,-0.809) (1.189,-0.828) (1.219,-0.845) (1.251,-0.861) (1.283,-0.876) (1.315,-0.889) (1.349,-0.900) (1.383,-0.909) (1.417,-0.917) (1.452,-0.923)};
\addplot[equipot] coordinates {(1.468,-0.655) (1.503,-0.658) (1.538,-0.657) (1.573,-0.653) (1.607,-0.645) (1.640,-0.634) (1.673,-0.620) (1.704,-0.604) (1.733,-0.585) (1.762,-0.564) (1.789,-0.542) (1.814,-0.518) (1.838,-0.492) (1.861,-0.465) (1.882,-0.437) (1.901,-0.407) (1.919,-0.377) (1.935,-0.346) (1.950,-0.314) (1.963,-0.282) (1.975,-0.249) (1.985,-0.215) (1.994,-0.181) (2.001,-0.146) (2.006,-0.112) (2.011,-0.077) (2.013,-0.042) (2.014,-0.007) (2.014,0.028) (2.012,0.063) (2.008,0.098) (2.003,0.133) (1.997,0.168) (1.989,0.202) (1.979,0.236) (1.968,0.269) (1.955,0.302) (1.941,0.334) (1.925,0.365) (1.908,0.396) (1.889,0.426) (1.869,0.454) (1.847,0.482) (1.824,0.508) (1.799,0.533) (1.772,0.556) (1.745,0.577) (1.715,0.597) (1.685,0.614) (1.653,0.629) (1.620,0.641) (1.586,0.650) (1.551,0.656) (1.516,0.658) (1.481,0.657) (1.446,0.652) (1.412,0.643) (1.379,0.632) (1.347,0.617) (1.317,0.600) (1.288,0.580) (1.260,0.558) (1.234,0.534) (1.210,0.509) (1.187,0.482) (1.166,0.454) (1.147,0.425) (1.129,0.395) (1.113,0.363) (1.098,0.331) (1.085,0.299) (1.073,0.266) (1.063,0.232) (1.054,0.198) (1.047,0.164) (1.041,0.129) (1.036,0.094) (1.033,0.059) (1.031,0.024) (1.031,-0.011) (1.032,-0.046) (1.035,-0.081) (1.039,-0.116) (1.044,-0.151) (1.051,-0.185) (1.059,-0.219) (1.069,-0.253) (1.080,-0.286) (1.093,-0.319) (1.107,-0.351) (1.123,-0.383) (1.140,-0.413) (1.159,-0.443) (1.179,-0.472) (1.201,-0.499) (1.225,-0.525) (1.250,-0.549) (1.277,-0.572) (1.305,-0.593) (1.335,-0.611) (1.367,-0.627) (1.399,-0.639) (1.433,-0.649) (1.468,-0.655)};
\addplot[equipot] coordinates {(1.468,-0.532) (1.503,-0.537) (1.537,-0.532) (1.570,-0.519) (1.600,-0.501) (1.629,-0.480) (1.654,-0.456) (1.678,-0.430) (1.700,-0.402) (1.720,-0.373) (1.738,-0.342) (1.754,-0.311) (1.769,-0.279) (1.782,-0.246) (1.793,-0.212) (1.802,-0.178) (1.809,-0.144) (1.815,-0.109) (1.820,-0.074) (1.822,-0.039) (1.823,-0.004) (1.823,0.032) (1.820,0.067) (1.816,0.102) (1.811,0.137) (1.804,0.171) (1.795,0.205) (1.784,0.239) (1.772,0.272) (1.757,0.304) (1.742,0.336) (1.724,0.367) (1.704,0.396) (1.683,0.424) (1.660,0.450) (1.634,0.475) (1.606,0.497) (1.577,0.515) (1.544,0.529) (1.510,0.537) (1.475,0.534) (1.442,0.522) (1.411,0.504) (1.383,0.483) (1.357,0.459) (1.334,0.432) (1.313,0.404) (1.293,0.374) (1.276,0.344) (1.260,0.312) (1.247,0.280) (1.234,0.246) (1.224,0.213) (1.215,0.179) (1.208,0.144) (1.202,0.109) (1.198,0.074) (1.196,0.039) (1.195,0.004) (1.196,-0.032) (1.198,-0.067) (1.201,-0.102) (1.207,-0.137) (1.214,-0.171) (1.222,-0.206) (1.232,-0.240) (1.244,-0.273) (1.258,-0.305) (1.273,-0.337) (1.290,-0.368) (1.309,-0.398) (1.330,-0.426) (1.352,-0.453) (1.378,-0.478) (1.405,-0.500) (1.435,-0.519) (1.468,-0.532)};
\addplot[equipot] coordinates {(1.500,-0.482) (1.526,-0.459) (1.550,-0.434) (1.572,-0.406) (1.592,-0.378) (1.610,-0.347) (1.627,-0.316) (1.641,-0.284) (1.654,-0.252) (1.665,-0.218) (1.674,-0.185) (1.682,-0.150) (1.688,-0.116) (1.692,-0.081) (1.695,-0.046) (1.697,-0.011) (1.696,0.024) (1.694,0.059) (1.691,0.094) (1.686,0.129) (1.679,0.163) (1.671,0.197) (1.661,0.231) (1.649,0.264) (1.636,0.297) (1.621,0.328) (1.604,0.359) (1.585,0.389) (1.564,0.417) (1.541,0.444) (1.516,0.468) (1.490,0.474) (1.465,0.449) (1.442,0.422) (1.422,0.394) (1.403,0.364) (1.386,0.333) (1.372,0.301) (1.358,0.269) (1.347,0.235) (1.337,0.202) (1.329,0.168) (1.323,0.133) (1.317,0.098) (1.314,0.063) (1.312,0.028) (1.312,-0.007) (1.313,-0.042) (1.315,-0.077) (1.319,-0.112) (1.325,-0.146) (1.332,-0.181) (1.341,-0.215) (1.351,-0.248) (1.363,-0.281) (1.377,-0.313) (1.392,-0.345) (1.410,-0.375) (1.429,-0.405) (1.451,-0.432) (1.474,-0.458) (1.500,-0.482)};
\addplot[equipot] coordinates {(1.500,-0.395) (1.516,-0.364) (1.531,-0.332) (1.545,-0.300) (1.557,-0.267) (1.568,-0.234) (1.577,-0.200) (1.585,-0.166) (1.591,-0.131) (1.596,-0.097) (1.599,-0.062) (1.601,-0.027) (1.601,0.008) (1.600,0.043) (1.598,0.078) (1.594,0.113) (1.588,0.147) (1.581,0.182) (1.573,0.216) (1.563,0.249) (1.552,0.282) (1.539,0.315) (1.524,0.347) (1.508,0.378) (1.494,0.380) (1.477,0.349) (1.463,0.317) (1.451,0.284) (1.439,0.251) (1.430,0.217) (1.422,0.183) (1.415,0.149) (1.410,0.114) (1.406,0.079) (1.403,0.044) (1.402,0.009) (1.402,-0.026) (1.404,-0.061) (1.407,-0.095) (1.412,-0.130) (1.418,-0.165) (1.425,-0.199) (1.434,-0.233) (1.444,-0.266) (1.456,-0.299) (1.470,-0.332) (1.484,-0.363) (1.500,-0.395)};
\addplot[lineaD, flechaen=0.169, flechaen=0.831] coordinates {(-3.200,0.617) (-3.167,0.604) (-3.134,0.592) (-3.102,0.580) (-3.069,0.567) (-3.036,0.555) (-3.003,0.543) (-2.970,0.531) (-2.937,0.518) (-2.904,0.506) (-2.871,0.494) (-2.838,0.483) (-2.805,0.471) (-2.772,0.459) (-2.739,0.447) (-2.706,0.436) (-2.673,0.424) (-2.640,0.413) (-2.607,0.401) (-2.574,0.390) (-2.541,0.379) (-2.507,0.367) (-2.474,0.356) (-2.441,0.345) (-2.407,0.334) (-2.374,0.324) (-2.341,0.313) (-2.307,0.302) (-2.274,0.292) (-2.240,0.282) (-2.207,0.271) (-2.173,0.261) (-2.140,0.251) (-2.106,0.241) (-2.072,0.232) (-2.038,0.222) (-2.005,0.213) (-1.971,0.204) (-1.937,0.195) (-1.903,0.186) (-1.869,0.178) (-1.835,0.170) (-1.800,0.162) (-1.766,0.155) (-1.732,0.148) (-1.697,0.141) (-1.663,0.135) (-1.628,0.129) (-1.594,0.123) (-1.559,0.118) (-1.524,0.114) (-1.489,0.110) (-1.455,0.106) (-1.420,0.103) (-1.385,0.100) (-1.350,0.098) (-1.315,0.095) (-1.280,0.093) (-1.245,0.092) (-1.210,0.090) (-1.175,0.089) (-1.140,0.088) (-1.104,0.087) (-1.069,0.086) (-1.034,0.085) (-0.999,0.085) (-0.964,0.084) (-0.929,0.084) (-0.894,0.083) (-0.859,0.083) (-0.824,0.083) (-0.789,0.082) (-0.754,0.082) (-0.719,0.082) (-0.684,0.081) (-0.649,0.081) (-0.614,0.081) (-0.579,0.081) (-0.544,0.081) (-0.508,0.081) (-0.473,0.081) (-0.438,0.080) (-0.403,0.080) (-0.368,0.080) (-0.333,0.080) (-0.298,0.080) (-0.263,0.080) (-0.228,0.080) (-0.193,0.080) (-0.158,0.080) (-0.123,0.080) (-0.088,0.080) (-0.053,0.080) (-0.018,0.080) (0.018,0.080) (0.053,0.080) (0.088,0.080) (0.123,0.080) (0.158,0.080) (0.193,0.080) (0.228,0.080) (0.263,0.080) (0.298,0.080) (0.333,0.080) (0.368,0.080) (0.403,0.080) (0.438,0.080) (0.473,0.081) (0.508,0.081) (0.544,0.081) (0.579,0.081) (0.614,0.081) (0.649,0.081) (0.684,0.081) (0.719,0.082) (0.754,0.082) (0.789,0.082) (0.824,0.083) (0.859,0.083) (0.894,0.083) (0.929,0.084) (0.964,0.084) (0.999,0.085) (1.034,0.085) (1.069,0.086) (1.104,0.087) (1.140,0.088) (1.175,0.089) (1.210,0.090) (1.245,0.092) (1.280,0.093) (1.315,0.095) (1.350,0.098) (1.385,0.100) (1.420,0.103) (1.455,0.106) (1.489,0.110) (1.524,0.114) (1.559,0.118) (1.594,0.123) (1.628,0.129) (1.663,0.135) (1.697,0.141) (1.732,0.148) (1.766,0.155) (1.800,0.162) (1.835,0.170) (1.869,0.178) (1.903,0.186) (1.937,0.195) (1.971,0.204) (2.005,0.213) (2.038,0.222) (2.072,0.232) (2.106,0.241) (2.140,0.251) (2.173,0.261) (2.207,0.271) (2.240,0.282) (2.274,0.292) (2.307,0.302) (2.341,0.313) (2.374,0.324) (2.407,0.334) (2.441,0.345) (2.474,0.356) (2.507,0.367) (2.541,0.379) (2.574,0.390) (2.607,0.401) (2.640,0.413) (2.673,0.424) (2.706,0.436) (2.739,0.447) (2.772,0.459) (2.805,0.471) (2.838,0.483) (2.871,0.494) (2.904,0.506) (2.937,0.518) (2.970,0.531) (3.003,0.543) (3.036,0.555) (3.069,0.567) (3.102,0.580) (3.134,0.592) (3.167,0.604) (3.200,0.617)};
\addplot[lineaD, flechaen=0.190, flechaen=0.810] coordinates {(-3.200,1.424) (-3.175,1.399) (-3.150,1.375) (-3.124,1.351) (-3.099,1.327) (-3.074,1.302) (-3.048,1.279) (-3.023,1.255) (-2.997,1.231) (-2.971,1.207) (-2.945,1.184) (-2.919,1.160) (-2.893,1.137) (-2.867,1.113) (-2.841,1.090) (-2.814,1.067) (-2.788,1.044) (-2.762,1.021) (-2.735,0.998) (-2.708,0.976) (-2.682,0.953) (-2.655,0.931) (-2.628,0.908) (-2.601,0.886) (-2.574,0.864) (-2.546,0.842) (-2.519,0.820) (-2.492,0.798) (-2.464,0.776) (-2.437,0.755) (-2.409,0.733) (-2.381,0.712) (-2.353,0.691) (-2.325,0.670) (-2.297,0.649) (-2.269,0.628) (-2.240,0.608) (-2.212,0.587) (-2.183,0.567) (-2.155,0.547) (-2.126,0.527) (-2.097,0.508) (-2.068,0.488) (-2.038,0.469) (-2.009,0.450) (-1.979,0.431) (-1.949,0.413) (-1.919,0.395) (-1.889,0.378) (-1.859,0.360) (-1.828,0.344) (-1.797,0.328) (-1.765,0.312) (-1.734,0.297) (-1.702,0.283) (-1.669,0.270) (-1.636,0.257) (-1.603,0.246) (-1.570,0.235) (-1.536,0.226) (-1.502,0.217) (-1.468,0.209) (-1.434,0.203) (-1.399,0.197) (-1.365,0.191) (-1.330,0.187) (-1.295,0.183) (-1.260,0.180) (-1.225,0.177) (-1.190,0.174) (-1.155,0.172) (-1.121,0.170) (-1.086,0.168) (-1.051,0.167) (-1.016,0.166) (-0.981,0.164) (-0.946,0.163) (-0.911,0.163) (-0.875,0.162) (-0.840,0.161) (-0.805,0.161) (-0.770,0.160) (-0.735,0.159) (-0.700,0.159) (-0.665,0.159) (-0.630,0.158) (-0.595,0.158) (-0.560,0.158) (-0.525,0.157) (-0.490,0.157) (-0.455,0.157) (-0.420,0.157) (-0.385,0.157) (-0.350,0.157) (-0.315,0.156) (-0.280,0.156) (-0.245,0.156) (-0.210,0.156) (-0.175,0.156) (-0.140,0.156) (-0.105,0.156) (-0.070,0.156) (-0.035,0.156) (0.000,0.156) (0.035,0.156) (0.070,0.156) (0.105,0.156) (0.140,0.156) (0.175,0.156) (0.210,0.156) (0.245,0.156) (0.280,0.156) (0.315,0.156) (0.350,0.157) (0.385,0.157) (0.420,0.157) (0.455,0.157) (0.490,0.157) (0.525,0.157) (0.560,0.158) (0.595,0.158) (0.630,0.158) (0.665,0.159) (0.700,0.159) (0.735,0.159) (0.770,0.160) (0.805,0.161) (0.840,0.161) (0.875,0.162) (0.911,0.163) (0.946,0.163) (0.981,0.164) (1.016,0.166) (1.051,0.167) (1.086,0.168) (1.121,0.170) (1.155,0.172) (1.190,0.174) (1.225,0.177) (1.260,0.180) (1.295,0.183) (1.330,0.187) (1.365,0.191) (1.399,0.197) (1.434,0.203) (1.468,0.209) (1.502,0.217) (1.536,0.226) (1.570,0.235) (1.603,0.246) (1.636,0.257) (1.669,0.270) (1.702,0.283) (1.734,0.297) (1.765,0.312) (1.797,0.328) (1.828,0.344) (1.859,0.360) (1.889,0.378) (1.919,0.395) (1.949,0.413) (1.979,0.431) (2.009,0.450) (2.038,0.469) (2.068,0.488) (2.097,0.508) (2.126,0.527) (2.155,0.547) (2.183,0.567) (2.212,0.587) (2.240,0.608) (2.269,0.628) (2.297,0.649) (2.325,0.670) (2.353,0.691) (2.381,0.712) (2.409,0.733) (2.437,0.755) (2.464,0.776) (2.492,0.798) (2.519,0.820) (2.546,0.842) (2.574,0.864) (2.601,0.886) (2.628,0.908) (2.655,0.931) (2.682,0.953) (2.708,0.976) (2.735,0.998) (2.762,1.021) (2.788,1.044) (2.814,1.067) (2.841,1.090) (2.867,1.113) (2.893,1.137) (2.919,1.160) (2.945,1.184) (2.971,1.207) (2.997,1.231) (3.023,1.255) (3.048,1.279) (3.074,1.302) (3.099,1.327) (3.124,1.351) (3.150,1.375) (3.175,1.399) (3.200,1.424)};
\addplot[lineaD, flechaen=0.138, flechaen=0.862] coordinates {(-2.570,1.900) (-2.557,1.867) (-2.544,1.835) (-2.531,1.802) (-2.518,1.770) (-2.504,1.738) (-2.490,1.705) (-2.477,1.673) (-2.462,1.641) (-2.448,1.609) (-2.434,1.577) (-2.419,1.546) (-2.404,1.514) (-2.389,1.483) (-2.373,1.451) (-2.358,1.420) (-2.342,1.388) (-2.326,1.357) (-2.310,1.326) (-2.293,1.295) (-2.277,1.264) (-2.260,1.234) (-2.243,1.203) (-2.226,1.173) (-2.208,1.142) (-2.191,1.112) (-2.173,1.082) (-2.155,1.052) (-2.137,1.022) (-2.118,0.992) (-2.100,0.962) (-2.081,0.933) (-2.062,0.904) (-2.042,0.874) (-2.023,0.845) (-2.003,0.816) (-1.983,0.788) (-1.963,0.759) (-1.942,0.731) (-1.921,0.703) (-1.900,0.675) (-1.879,0.647) (-1.857,0.620) (-1.835,0.593) (-1.812,0.566) (-1.789,0.540) (-1.765,0.514) (-1.741,0.489) (-1.716,0.464) (-1.690,0.441) (-1.663,0.418) (-1.635,0.397) (-1.607,0.377) (-1.577,0.359) (-1.546,0.342) (-1.514,0.328) (-1.481,0.315) (-1.448,0.304) (-1.414,0.294) (-1.380,0.286) (-1.346,0.279) (-1.312,0.273) (-1.277,0.268) (-1.242,0.263) (-1.207,0.259) (-1.173,0.256) (-1.138,0.253) (-1.103,0.251) (-1.068,0.248) (-1.033,0.247) (-0.998,0.245) (-0.963,0.243) (-0.928,0.242) (-0.893,0.241) (-0.858,0.240) (-0.823,0.239) (-0.788,0.238) (-0.753,0.237) (-0.718,0.237) (-0.683,0.236) (-0.648,0.236) (-0.613,0.235) (-0.578,0.235) (-0.543,0.234) (-0.508,0.234) (-0.473,0.234) (-0.438,0.233) (-0.403,0.233) (-0.368,0.233) (-0.333,0.233) (-0.298,0.233) (-0.263,0.232) (-0.228,0.232) (-0.193,0.232) (-0.158,0.232) (-0.123,0.232) (-0.088,0.232) (-0.053,0.232) (-0.018,0.232) (0.018,0.232) (0.053,0.232) (0.088,0.232) (0.123,0.232) (0.158,0.232) (0.193,0.232) (0.228,0.232) (0.263,0.232) (0.298,0.233) (0.333,0.233) (0.368,0.233) (0.403,0.233) (0.438,0.233) (0.473,0.234) (0.508,0.234) (0.543,0.234) (0.578,0.235) (0.613,0.235) (0.648,0.236) (0.683,0.236) (0.718,0.237) (0.753,0.237) (0.788,0.238) (0.823,0.239) (0.858,0.240) (0.893,0.241) (0.928,0.242) (0.963,0.243) (0.998,0.245) (1.033,0.247) (1.068,0.248) (1.103,0.251) (1.138,0.253) (1.173,0.256) (1.207,0.259) (1.242,0.263) (1.277,0.268) (1.312,0.273) (1.346,0.279) (1.380,0.286) (1.414,0.294) (1.448,0.304) (1.481,0.315) (1.514,0.328) (1.546,0.342) (1.577,0.359) (1.607,0.377) (1.635,0.397) (1.663,0.418) (1.690,0.441) (1.716,0.464) (1.741,0.489) (1.765,0.514) (1.789,0.540) (1.812,0.566) (1.835,0.593) (1.857,0.620) (1.879,0.647) (1.900,0.675) (1.921,0.703) (1.942,0.731) (1.963,0.759) (1.983,0.788) (2.003,0.816) (2.023,0.845) (2.042,0.874) (2.062,0.904) (2.081,0.933) (2.100,0.962) (2.118,0.992) (2.137,1.022) (2.155,1.052) (2.173,1.082) (2.191,1.112) (2.208,1.142) (2.226,1.173) (2.243,1.203) (2.260,1.234) (2.277,1.264) (2.293,1.295) (2.310,1.326) (2.326,1.357) (2.342,1.388) (2.358,1.420) (2.373,1.451) (2.389,1.483) (2.404,1.514) (2.419,1.546) (2.434,1.577) (2.448,1.609) (2.462,1.641) (2.477,1.673) (2.490,1.705) (2.504,1.738) (2.518,1.770) (2.531,1.802) (2.544,1.835) (2.557,1.867) (2.570,1.900)};
\addplot[lineaD, flechaen=0.582] coordinates {(-1.772,1.900) (-1.774,1.865) (-1.776,1.830) (-1.778,1.795) (-1.779,1.759) (-1.780,1.724) (-1.781,1.689) (-1.781,1.654) (-1.781,1.619) (-1.780,1.583) (-1.779,1.548) (-1.778,1.513) (-1.776,1.478) (-1.774,1.443) (-1.772,1.408) (-1.769,1.373) (-1.766,1.338) (-1.763,1.303) (-1.759,1.268) (-1.755,1.233) (-1.750,1.198) (-1.746,1.163) (-1.741,1.128) (-1.735,1.093) (-1.729,1.059) (-1.723,1.024) (-1.717,0.989) (-1.710,0.955) (-1.703,0.920) (-1.695,0.886) (-1.687,0.852) (-1.679,0.817) (-1.670,0.783) (-1.661,0.749) (-1.652,0.715) (-1.641,0.682) (-1.631,0.648) (-1.619,0.615) (-1.607,0.582) (-1.593,0.550) (-1.578,0.518) (-1.560,0.488) (-1.538,0.460) (-1.512,0.436) (-1.483,0.417) (-1.452,0.401) (-1.419,0.388) (-1.386,0.377) (-1.352,0.368) (-1.317,0.360) (-1.283,0.353) (-1.248,0.348) (-1.213,0.343) (-1.178,0.338) (-1.143,0.335) (-1.108,0.332) (-1.073,0.329) (-1.038,0.326) (-1.003,0.324) (-0.968,0.322) (-0.933,0.321) (-0.897,0.319) (-0.862,0.318) (-0.827,0.317) (-0.792,0.316) (-0.757,0.315) (-0.721,0.314) (-0.686,0.313) (-0.651,0.313) (-0.616,0.312) (-0.581,0.311) (-0.546,0.311) (-0.510,0.311) (-0.475,0.310) (-0.440,0.310) (-0.405,0.310) (-0.370,0.309) (-0.334,0.309) (-0.299,0.309) (-0.264,0.309) (-0.229,0.308) (-0.194,0.308) (-0.158,0.308) (-0.123,0.308) (-0.088,0.308) (-0.053,0.308) (-0.018,0.308) (0.018,0.308) (0.053,0.308) (0.088,0.308) (0.123,0.308) (0.158,0.308) (0.194,0.308) (0.229,0.308) (0.264,0.309) (0.299,0.309) (0.334,0.309) (0.370,0.309) (0.405,0.310) (0.440,0.310) (0.475,0.310) (0.510,0.311) (0.546,0.311) (0.581,0.311) (0.616,0.312) (0.651,0.313) (0.686,0.313) (0.721,0.314) (0.757,0.315) (0.792,0.316) (0.827,0.317) (0.862,0.318) (0.897,0.319) (0.933,0.321) (0.968,0.322) (1.003,0.324) (1.038,0.326) (1.073,0.329) (1.108,0.332) (1.143,0.335) (1.178,0.338) (1.213,0.343) (1.248,0.348) (1.283,0.353) (1.317,0.360) (1.352,0.368) (1.386,0.377) (1.419,0.388) (1.452,0.401) (1.483,0.417) (1.512,0.436) (1.538,0.460) (1.560,0.488) (1.578,0.518) (1.593,0.550) (1.607,0.582) (1.619,0.615) (1.631,0.648) (1.641,0.682) (1.652,0.715) (1.661,0.749) (1.670,0.783) (1.679,0.817) (1.687,0.852) (1.695,0.886) (1.703,0.920) (1.710,0.955) (1.717,0.989) (1.723,1.024) (1.729,1.059) (1.735,1.093) (1.741,1.128) (1.746,1.163) (1.750,1.198) (1.755,1.233) (1.759,1.268) (1.763,1.303) (1.766,1.338) (1.769,1.373) (1.772,1.408) (1.774,1.443) (1.776,1.478) (1.778,1.513) (1.779,1.548) (1.780,1.583) (1.781,1.619) (1.781,1.654) (1.781,1.689) (1.780,1.724) (1.779,1.759) (1.778,1.795) (1.776,1.830) (1.774,1.865) (1.772,1.900)};
\addplot[lineaD, flechaen=0.418] coordinates {(-0.818,1.900) (-0.843,1.876) (-0.868,1.851) (-0.892,1.825) (-0.916,1.799) (-0.939,1.772) (-0.961,1.745) (-0.983,1.718) (-1.004,1.690) (-1.024,1.661) (-1.044,1.632) (-1.064,1.603) (-1.083,1.574) (-1.101,1.544) (-1.119,1.513) (-1.136,1.483) (-1.153,1.452) (-1.169,1.421) (-1.185,1.389) (-1.200,1.358) (-1.215,1.326) (-1.229,1.294) (-1.243,1.262) (-1.256,1.229) (-1.269,1.196) (-1.282,1.164) (-1.294,1.131) (-1.306,1.098) (-1.317,1.064) (-1.328,1.031) (-1.338,0.997) (-1.348,0.964) (-1.358,0.930) (-1.368,0.896) (-1.377,0.862) (-1.386,0.828) (-1.395,0.794) (-1.403,0.760) (-1.411,0.726) (-1.420,0.692) (-1.428,0.658) (-1.436,0.624) (-1.445,0.590) (-1.454,0.556) (-1.464,0.522) (-1.462,0.493) (-1.431,0.478) (-1.398,0.465) (-1.365,0.454) (-1.331,0.445) (-1.297,0.437) (-1.262,0.431) (-1.228,0.425) (-1.193,0.420) (-1.158,0.416) (-1.123,0.412) (-1.088,0.409) (-1.053,0.406) (-1.018,0.404) (-0.983,0.402) (-0.948,0.400) (-0.913,0.398) (-0.878,0.396) (-0.843,0.395) (-0.808,0.394) (-0.772,0.393) (-0.737,0.392) (-0.702,0.391) (-0.667,0.390) (-0.632,0.389) (-0.597,0.388) (-0.562,0.388) (-0.527,0.387) (-0.492,0.387) (-0.456,0.386) (-0.421,0.386) (-0.386,0.386) (-0.351,0.385) (-0.316,0.385) (-0.281,0.385) (-0.246,0.385) (-0.211,0.384) (-0.176,0.384) (-0.140,0.384) (-0.105,0.384) (-0.070,0.384) (-0.035,0.384) (0.000,0.384) (0.035,0.384) (0.070,0.384) (0.105,0.384) (0.140,0.384) (0.176,0.384) (0.211,0.384) (0.246,0.385) (0.281,0.385) (0.316,0.385) (0.351,0.385) (0.386,0.386) (0.421,0.386) (0.456,0.386) (0.492,0.387) (0.527,0.387) (0.562,0.388) (0.597,0.388) (0.632,0.389) (0.667,0.390) (0.702,0.391) (0.737,0.392) (0.772,0.393) (0.808,0.394) (0.843,0.395) (0.878,0.396) (0.913,0.398) (0.948,0.400) (0.983,0.402) (1.018,0.404) (1.053,0.406) (1.088,0.409) (1.123,0.412) (1.158,0.416) (1.193,0.420) (1.228,0.425) (1.262,0.431) (1.297,0.437) (1.331,0.445) (1.365,0.454) (1.398,0.465) (1.431,0.478) (1.462,0.493) (1.464,0.522) (1.454,0.556) (1.445,0.590) (1.436,0.624) (1.428,0.658) (1.420,0.692) (1.411,0.726) (1.403,0.760) (1.395,0.794) (1.386,0.828) (1.377,0.862) (1.368,0.896) (1.358,0.930) (1.348,0.964) (1.338,0.997) (1.328,1.031) (1.317,1.064) (1.306,1.098) (1.294,1.131) (1.282,1.164) (1.269,1.196) (1.256,1.229) (1.243,1.262) (1.229,1.294) (1.215,1.326) (1.200,1.358) (1.185,1.389) (1.169,1.421) (1.153,1.452) (1.136,1.483) (1.119,1.513) (1.101,1.544) (1.083,1.574) (1.064,1.603) (1.044,1.632) (1.024,1.661) (1.004,1.690) (0.983,1.718) (0.961,1.745) (0.939,1.772) (0.916,1.799) (0.892,1.825) (0.868,1.851) (0.843,1.876) (0.818,1.900)};
\addplot[lineaD, flechaen=0.209] coordinates {(-0.576,0.465) (-0.541,0.464) (-0.506,0.463) (-0.470,0.463) (-0.435,0.462) (-0.400,0.462) (-0.365,0.462) (-0.330,0.461) (-0.294,0.461) (-0.259,0.461) (-0.224,0.461) (-0.189,0.460) (-0.154,0.460) (-0.118,0.460) (-0.083,0.460) (-0.048,0.460) (-0.013,0.460) (0.022,0.460) (0.058,0.460) (0.093,0.460) (0.128,0.460) (0.163,0.460) (0.198,0.460) (0.234,0.461) (0.269,0.461) (0.304,0.461) (0.339,0.461) (0.374,0.462) (0.410,0.462) (0.445,0.463) (0.480,0.463) (0.515,0.464) (0.550,0.464) (0.586,0.465) (0.621,0.465) (0.656,0.466) (0.691,0.467) (0.726,0.468) (0.762,0.469) (0.797,0.470) (0.832,0.472) (0.867,0.473) (0.902,0.475) (0.937,0.477) (0.973,0.479) (1.008,0.481) (1.043,0.483) (1.078,0.486) (1.113,0.490) (1.148,0.493) (1.182,0.498) (1.183,0.517) (1.156,0.540) (1.130,0.564) (1.104,0.587) (1.077,0.611) (1.051,0.634) (1.025,0.657) (0.998,0.680) (0.971,0.703) (0.945,0.726) (0.918,0.749) (0.890,0.771) (0.863,0.793) (0.835,0.814) (0.807,0.835) (0.778,0.856) (0.749,0.876) (0.720,0.896) (0.691,0.915) (0.661,0.934) (0.630,0.952) (0.600,0.969) (0.569,0.986) (0.537,1.002) (0.506,1.017) (0.474,1.031) (0.441,1.045) (0.408,1.058) (0.375,1.069) (0.342,1.080) (0.308,1.090) (0.274,1.099) (0.240,1.107) (0.205,1.114) (0.170,1.120) (0.136,1.125) (0.101,1.129) (0.066,1.131) (0.030,1.133) (-0.005,1.133) (-0.040,1.133) (-0.075,1.131) (-0.110,1.128) (-0.145,1.124) (-0.180,1.119) (-0.215,1.113) (-0.249,1.105) (-0.283,1.097) (-0.317,1.088) (-0.351,1.077) (-0.384,1.066) (-0.417,1.054) (-0.450,1.041) (-0.482,1.027) (-0.514,1.013) (-0.546,0.997) (-0.577,0.981) (-0.608,0.964) (-0.639,0.947) (-0.669,0.928) (-0.699,0.910) (-0.728,0.890) (-0.757,0.871) (-0.786,0.850) (-0.814,0.829) (-0.842,0.808) (-0.870,0.787) (-0.898,0.765) (-0.925,0.742) (-0.952,0.720) (-0.979,0.697) (-1.005,0.674) (-1.032,0.651) (-1.058,0.628) (-1.085,0.604) (-1.111,0.581) (-1.137,0.557) (-1.164,0.534) (-1.190,0.511) (-1.173,0.496) (-1.138,0.492) (-1.103,0.489) (-1.068,0.486) (-1.033,0.483) (-0.998,0.480) (-0.963,0.478) (-0.928,0.476) (-0.893,0.474) (-0.857,0.473) (-0.822,0.471) (-0.787,0.470) (-0.752,0.469) (-0.717,0.468) (-0.681,0.467) (-0.646,0.466) (-0.611,0.465) (-0.576,0.465)};
\addplot[lineaD, flechaen=0.169, flechaen=0.831] coordinates {(-3.200,-0.617) (-3.167,-0.604) (-3.134,-0.592) (-3.102,-0.580) (-3.069,-0.567) (-3.036,-0.555) (-3.003,-0.543) (-2.970,-0.531) (-2.937,-0.518) (-2.904,-0.506) (-2.871,-0.494) (-2.838,-0.483) (-2.805,-0.471) (-2.772,-0.459) (-2.739,-0.447) (-2.706,-0.436) (-2.673,-0.424) (-2.640,-0.413) (-2.607,-0.401) (-2.574,-0.390) (-2.541,-0.379) (-2.507,-0.367) (-2.474,-0.356) (-2.441,-0.345) (-2.407,-0.334) (-2.374,-0.324) (-2.341,-0.313) (-2.307,-0.302) (-2.274,-0.292) (-2.240,-0.282) (-2.207,-0.271) (-2.173,-0.261) (-2.140,-0.251) (-2.106,-0.241) (-2.072,-0.232) (-2.038,-0.222) (-2.005,-0.213) (-1.971,-0.204) (-1.937,-0.195) (-1.903,-0.186) (-1.869,-0.178) (-1.835,-0.170) (-1.800,-0.162) (-1.766,-0.155) (-1.732,-0.148) (-1.697,-0.141) (-1.663,-0.135) (-1.628,-0.129) (-1.594,-0.123) (-1.559,-0.118) (-1.524,-0.114) (-1.489,-0.110) (-1.455,-0.106) (-1.420,-0.103) (-1.385,-0.100) (-1.350,-0.098) (-1.315,-0.095) (-1.280,-0.093) (-1.245,-0.092) (-1.210,-0.090) (-1.175,-0.089) (-1.140,-0.088) (-1.104,-0.087) (-1.069,-0.086) (-1.034,-0.085) (-0.999,-0.085) (-0.964,-0.084) (-0.929,-0.084) (-0.894,-0.083) (-0.859,-0.083) (-0.824,-0.083) (-0.789,-0.082) (-0.754,-0.082) (-0.719,-0.082) (-0.684,-0.081) (-0.649,-0.081) (-0.614,-0.081) (-0.579,-0.081) (-0.544,-0.081) (-0.508,-0.081) (-0.473,-0.081) (-0.438,-0.080) (-0.403,-0.080) (-0.368,-0.080) (-0.333,-0.080) (-0.298,-0.080) (-0.263,-0.080) (-0.228,-0.080) (-0.193,-0.080) (-0.158,-0.080) (-0.123,-0.080) (-0.088,-0.080) (-0.053,-0.080) (-0.018,-0.080) (0.018,-0.080) (0.053,-0.080) (0.088,-0.080) (0.123,-0.080) (0.158,-0.080) (0.193,-0.080) (0.228,-0.080) (0.263,-0.080) (0.298,-0.080) (0.333,-0.080) (0.368,-0.080) (0.403,-0.080) (0.438,-0.080) (0.473,-0.081) (0.508,-0.081) (0.544,-0.081) (0.579,-0.081) (0.614,-0.081) (0.649,-0.081) (0.684,-0.081) (0.719,-0.082) (0.754,-0.082) (0.789,-0.082) (0.824,-0.083) (0.859,-0.083) (0.894,-0.083) (0.929,-0.084) (0.964,-0.084) (0.999,-0.085) (1.034,-0.085) (1.069,-0.086) (1.104,-0.087) (1.140,-0.088) (1.175,-0.089) (1.210,-0.090) (1.245,-0.092) (1.280,-0.093) (1.315,-0.095) (1.350,-0.098) (1.385,-0.100) (1.420,-0.103) (1.455,-0.106) (1.489,-0.110) (1.524,-0.114) (1.559,-0.118) (1.594,-0.123) (1.628,-0.129) (1.663,-0.135) (1.697,-0.141) (1.732,-0.148) (1.766,-0.155) (1.800,-0.162) (1.835,-0.170) (1.869,-0.178) (1.903,-0.186) (1.937,-0.195) (1.971,-0.204) (2.005,-0.213) (2.038,-0.222) (2.072,-0.232) (2.106,-0.241) (2.140,-0.251) (2.173,-0.261) (2.207,-0.271) (2.240,-0.282) (2.274,-0.292) (2.307,-0.302) (2.341,-0.313) (2.374,-0.324) (2.407,-0.334) (2.441,-0.345) (2.474,-0.356) (2.507,-0.367) (2.541,-0.379) (2.574,-0.390) (2.607,-0.401) (2.640,-0.413) (2.673,-0.424) (2.706,-0.436) (2.739,-0.447) (2.772,-0.459) (2.805,-0.471) (2.838,-0.483) (2.871,-0.494) (2.904,-0.506) (2.937,-0.518) (2.970,-0.531) (3.003,-0.543) (3.036,-0.555) (3.069,-0.567) (3.102,-0.580) (3.134,-0.592) (3.167,-0.604) (3.200,-0.617)};
\addplot[lineaD, flechaen=0.190, flechaen=0.810] coordinates {(-3.200,-1.424) (-3.175,-1.399) (-3.150,-1.375) (-3.124,-1.351) (-3.099,-1.327) (-3.074,-1.302) (-3.048,-1.279) (-3.023,-1.255) (-2.997,-1.231) (-2.971,-1.207) (-2.945,-1.184) (-2.919,-1.160) (-2.893,-1.137) (-2.867,-1.113) (-2.841,-1.090) (-2.814,-1.067) (-2.788,-1.044) (-2.762,-1.021) (-2.735,-0.998) (-2.708,-0.976) (-2.682,-0.953) (-2.655,-0.931) (-2.628,-0.908) (-2.601,-0.886) (-2.574,-0.864) (-2.546,-0.842) (-2.519,-0.820) (-2.492,-0.798) (-2.464,-0.776) (-2.437,-0.755) (-2.409,-0.733) (-2.381,-0.712) (-2.353,-0.691) (-2.325,-0.670) (-2.297,-0.649) (-2.269,-0.628) (-2.240,-0.608) (-2.212,-0.587) (-2.183,-0.567) (-2.155,-0.547) (-2.126,-0.527) (-2.097,-0.508) (-2.068,-0.488) (-2.038,-0.469) (-2.009,-0.450) (-1.979,-0.431) (-1.949,-0.413) (-1.919,-0.395) (-1.889,-0.378) (-1.859,-0.360) (-1.828,-0.344) (-1.797,-0.328) (-1.765,-0.312) (-1.734,-0.297) (-1.702,-0.283) (-1.669,-0.270) (-1.636,-0.257) (-1.603,-0.246) (-1.570,-0.235) (-1.536,-0.226) (-1.502,-0.217) (-1.468,-0.209) (-1.434,-0.203) (-1.399,-0.197) (-1.365,-0.191) (-1.330,-0.187) (-1.295,-0.183) (-1.260,-0.180) (-1.225,-0.177) (-1.190,-0.174) (-1.155,-0.172) (-1.121,-0.170) (-1.086,-0.168) (-1.051,-0.167) (-1.016,-0.166) (-0.981,-0.164) (-0.946,-0.163) (-0.911,-0.163) (-0.875,-0.162) (-0.840,-0.161) (-0.805,-0.161) (-0.770,-0.160) (-0.735,-0.159) (-0.700,-0.159) (-0.665,-0.159) (-0.630,-0.158) (-0.595,-0.158) (-0.560,-0.158) (-0.525,-0.157) (-0.490,-0.157) (-0.455,-0.157) (-0.420,-0.157) (-0.385,-0.157) (-0.350,-0.157) (-0.315,-0.156) (-0.280,-0.156) (-0.245,-0.156) (-0.210,-0.156) (-0.175,-0.156) (-0.140,-0.156) (-0.105,-0.156) (-0.070,-0.156) (-0.035,-0.156) (0.000,-0.156) (0.035,-0.156) (0.070,-0.156) (0.105,-0.156) (0.140,-0.156) (0.175,-0.156) (0.210,-0.156) (0.245,-0.156) (0.280,-0.156) (0.315,-0.156) (0.350,-0.157) (0.385,-0.157) (0.420,-0.157) (0.455,-0.157) (0.490,-0.157) (0.525,-0.157) (0.560,-0.158) (0.595,-0.158) (0.630,-0.158) (0.665,-0.159) (0.700,-0.159) (0.735,-0.159) (0.770,-0.160) (0.805,-0.161) (0.840,-0.161) (0.875,-0.162) (0.911,-0.163) (0.946,-0.163) (0.981,-0.164) (1.016,-0.166) (1.051,-0.167) (1.086,-0.168) (1.121,-0.170) (1.155,-0.172) (1.190,-0.174) (1.225,-0.177) (1.260,-0.180) (1.295,-0.183) (1.330,-0.187) (1.365,-0.191) (1.399,-0.197) (1.434,-0.203) (1.468,-0.209) (1.502,-0.217) (1.536,-0.226) (1.570,-0.235) (1.603,-0.246) (1.636,-0.257) (1.669,-0.270) (1.702,-0.283) (1.734,-0.297) (1.765,-0.312) (1.797,-0.328) (1.828,-0.344) (1.859,-0.360) (1.889,-0.378) (1.919,-0.395) (1.949,-0.413) (1.979,-0.431) (2.009,-0.450) (2.038,-0.469) (2.068,-0.488) (2.097,-0.508) (2.126,-0.527) (2.155,-0.547) (2.183,-0.567) (2.212,-0.587) (2.240,-0.608) (2.269,-0.628) (2.297,-0.649) (2.325,-0.670) (2.353,-0.691) (2.381,-0.712) (2.409,-0.733) (2.437,-0.755) (2.464,-0.776) (2.492,-0.798) (2.519,-0.820) (2.546,-0.842) (2.574,-0.864) (2.601,-0.886) (2.628,-0.908) (2.655,-0.931) (2.682,-0.953) (2.708,-0.976) (2.735,-0.998) (2.762,-1.021) (2.788,-1.044) (2.814,-1.067) (2.841,-1.090) (2.867,-1.113) (2.893,-1.137) (2.919,-1.160) (2.945,-1.184) (2.971,-1.207) (2.997,-1.231) (3.023,-1.255) (3.048,-1.279) (3.074,-1.302) (3.099,-1.327) (3.124,-1.351) (3.150,-1.375) (3.175,-1.399) (3.200,-1.424)};
\addplot[lineaD, flechaen=0.138, flechaen=0.862] coordinates {(-2.570,-1.900) (-2.557,-1.867) (-2.544,-1.835) (-2.531,-1.802) (-2.518,-1.770) (-2.504,-1.738) (-2.490,-1.705) (-2.477,-1.673) (-2.462,-1.641) (-2.448,-1.609) (-2.434,-1.577) (-2.419,-1.546) (-2.404,-1.514) (-2.389,-1.483) (-2.373,-1.451) (-2.358,-1.420) (-2.342,-1.388) (-2.326,-1.357) (-2.310,-1.326) (-2.293,-1.295) (-2.277,-1.264) (-2.260,-1.234) (-2.243,-1.203) (-2.226,-1.173) (-2.208,-1.142) (-2.191,-1.112) (-2.173,-1.082) (-2.155,-1.052) (-2.137,-1.022) (-2.118,-0.992) (-2.100,-0.962) (-2.081,-0.933) (-2.062,-0.904) (-2.042,-0.874) (-2.023,-0.845) (-2.003,-0.816) (-1.983,-0.788) (-1.963,-0.759) (-1.942,-0.731) (-1.921,-0.703) (-1.900,-0.675) (-1.879,-0.647) (-1.857,-0.620) (-1.835,-0.593) (-1.812,-0.566) (-1.789,-0.540) (-1.765,-0.514) (-1.741,-0.489) (-1.716,-0.464) (-1.690,-0.441) (-1.663,-0.418) (-1.635,-0.397) (-1.607,-0.377) (-1.577,-0.359) (-1.546,-0.342) (-1.514,-0.328) (-1.481,-0.315) (-1.448,-0.304) (-1.414,-0.294) (-1.380,-0.286) (-1.346,-0.279) (-1.312,-0.273) (-1.277,-0.268) (-1.242,-0.263) (-1.207,-0.259) (-1.173,-0.256) (-1.138,-0.253) (-1.103,-0.251) (-1.068,-0.248) (-1.033,-0.247) (-0.998,-0.245) (-0.963,-0.243) (-0.928,-0.242) (-0.893,-0.241) (-0.858,-0.240) (-0.823,-0.239) (-0.788,-0.238) (-0.753,-0.237) (-0.718,-0.237) (-0.683,-0.236) (-0.648,-0.236) (-0.613,-0.235) (-0.578,-0.235) (-0.543,-0.234) (-0.508,-0.234) (-0.473,-0.234) (-0.438,-0.233) (-0.403,-0.233) (-0.368,-0.233) (-0.333,-0.233) (-0.298,-0.233) (-0.263,-0.232) (-0.228,-0.232) (-0.193,-0.232) (-0.158,-0.232) (-0.123,-0.232) (-0.088,-0.232) (-0.053,-0.232) (-0.018,-0.232) (0.018,-0.232) (0.053,-0.232) (0.088,-0.232) (0.123,-0.232) (0.158,-0.232) (0.193,-0.232) (0.228,-0.232) (0.263,-0.232) (0.298,-0.233) (0.333,-0.233) (0.368,-0.233) (0.403,-0.233) (0.438,-0.233) (0.473,-0.234) (0.508,-0.234) (0.543,-0.234) (0.578,-0.235) (0.613,-0.235) (0.648,-0.236) (0.683,-0.236) (0.718,-0.237) (0.753,-0.237) (0.788,-0.238) (0.823,-0.239) (0.858,-0.240) (0.893,-0.241) (0.928,-0.242) (0.963,-0.243) (0.998,-0.245) (1.033,-0.247) (1.068,-0.248) (1.103,-0.251) (1.138,-0.253) (1.173,-0.256) (1.207,-0.259) (1.242,-0.263) (1.277,-0.268) (1.312,-0.273) (1.346,-0.279) (1.380,-0.286) (1.414,-0.294) (1.448,-0.304) (1.481,-0.315) (1.514,-0.328) (1.546,-0.342) (1.577,-0.359) (1.607,-0.377) (1.635,-0.397) (1.663,-0.418) (1.690,-0.441) (1.716,-0.464) (1.741,-0.489) (1.765,-0.514) (1.789,-0.540) (1.812,-0.566) (1.835,-0.593) (1.857,-0.620) (1.879,-0.647) (1.900,-0.675) (1.921,-0.703) (1.942,-0.731) (1.963,-0.759) (1.983,-0.788) (2.003,-0.816) (2.023,-0.845) (2.042,-0.874) (2.062,-0.904) (2.081,-0.933) (2.100,-0.962) (2.118,-0.992) (2.137,-1.022) (2.155,-1.052) (2.173,-1.082) (2.191,-1.112) (2.208,-1.142) (2.226,-1.173) (2.243,-1.203) (2.260,-1.234) (2.277,-1.264) (2.293,-1.295) (2.310,-1.326) (2.326,-1.357) (2.342,-1.388) (2.358,-1.420) (2.373,-1.451) (2.389,-1.483) (2.404,-1.514) (2.419,-1.546) (2.434,-1.577) (2.448,-1.609) (2.462,-1.641) (2.477,-1.673) (2.490,-1.705) (2.504,-1.738) (2.518,-1.770) (2.531,-1.802) (2.544,-1.835) (2.557,-1.867) (2.570,-1.900)};
\addplot[lineaD, flechaen=0.418] coordinates {(-1.772,-1.900) (-1.774,-1.865) (-1.776,-1.830) (-1.778,-1.795) (-1.779,-1.759) (-1.780,-1.724) (-1.781,-1.689) (-1.781,-1.654) (-1.781,-1.619) (-1.780,-1.583) (-1.779,-1.548) (-1.778,-1.513) (-1.776,-1.478) (-1.774,-1.443) (-1.772,-1.408) (-1.769,-1.373) (-1.766,-1.338) (-1.763,-1.303) (-1.759,-1.268) (-1.755,-1.233) (-1.750,-1.198) (-1.746,-1.163) (-1.741,-1.128) (-1.735,-1.093) (-1.729,-1.059) (-1.723,-1.024) (-1.717,-0.989) (-1.710,-0.955) (-1.703,-0.920) (-1.695,-0.886) (-1.687,-0.852) (-1.679,-0.817) (-1.670,-0.783) (-1.661,-0.749) (-1.652,-0.715) (-1.641,-0.682) (-1.631,-0.648) (-1.619,-0.615) (-1.607,-0.582) (-1.593,-0.550) (-1.578,-0.518) (-1.560,-0.488) (-1.538,-0.460) (-1.512,-0.436) (-1.483,-0.417) (-1.452,-0.401) (-1.419,-0.388) (-1.386,-0.377) (-1.352,-0.368) (-1.317,-0.360) (-1.283,-0.353) (-1.248,-0.348) (-1.213,-0.343) (-1.178,-0.338) (-1.143,-0.335) (-1.108,-0.332) (-1.073,-0.329) (-1.038,-0.326) (-1.003,-0.324) (-0.968,-0.322) (-0.933,-0.321) (-0.897,-0.319) (-0.862,-0.318) (-0.827,-0.317) (-0.792,-0.316) (-0.757,-0.315) (-0.721,-0.314) (-0.686,-0.313) (-0.651,-0.313) (-0.616,-0.312) (-0.581,-0.311) (-0.546,-0.311) (-0.510,-0.311) (-0.475,-0.310) (-0.440,-0.310) (-0.405,-0.310) (-0.370,-0.309) (-0.334,-0.309) (-0.299,-0.309) (-0.264,-0.309) (-0.229,-0.308) (-0.194,-0.308) (-0.158,-0.308) (-0.123,-0.308) (-0.088,-0.308) (-0.053,-0.308) (-0.018,-0.308) (0.018,-0.308) (0.053,-0.308) (0.088,-0.308) (0.123,-0.308) (0.158,-0.308) (0.194,-0.308) (0.229,-0.308) (0.264,-0.309) (0.299,-0.309) (0.334,-0.309) (0.370,-0.309) (0.405,-0.310) (0.440,-0.310) (0.475,-0.310) (0.510,-0.311) (0.546,-0.311) (0.581,-0.311) (0.616,-0.312) (0.651,-0.313) (0.686,-0.313) (0.721,-0.314) (0.757,-0.315) (0.792,-0.316) (0.827,-0.317) (0.862,-0.318) (0.897,-0.319) (0.933,-0.321) (0.968,-0.322) (1.003,-0.324) (1.038,-0.326) (1.073,-0.329) (1.108,-0.332) (1.143,-0.335) (1.178,-0.338) (1.213,-0.343) (1.248,-0.348) (1.283,-0.353) (1.317,-0.360) (1.352,-0.368) (1.386,-0.377) (1.419,-0.388) (1.452,-0.401) (1.483,-0.417) (1.512,-0.436) (1.538,-0.460) (1.560,-0.488) (1.578,-0.518) (1.593,-0.550) (1.607,-0.582) (1.619,-0.615) (1.631,-0.648) (1.641,-0.682) (1.652,-0.715) (1.661,-0.749) (1.670,-0.783) (1.679,-0.817) (1.687,-0.852) (1.695,-0.886) (1.703,-0.920) (1.710,-0.955) (1.717,-0.989) (1.723,-1.024) (1.729,-1.059) (1.735,-1.093) (1.741,-1.128) (1.746,-1.163) (1.750,-1.198) (1.755,-1.233) (1.759,-1.268) (1.763,-1.303) (1.766,-1.338) (1.769,-1.373) (1.772,-1.408) (1.774,-1.443) (1.776,-1.478) (1.778,-1.513) (1.779,-1.548) (1.780,-1.583) (1.781,-1.619) (1.781,-1.654) (1.781,-1.689) (1.780,-1.724) (1.779,-1.759) (1.778,-1.795) (1.776,-1.830) (1.774,-1.865) (1.772,-1.900)};
\addplot[lineaD, flechaen=0.582] coordinates {(-0.818,-1.900) (-0.843,-1.876) (-0.868,-1.851) (-0.892,-1.825) (-0.916,-1.799) (-0.939,-1.772) (-0.961,-1.745) (-0.983,-1.718) (-1.004,-1.690) (-1.024,-1.661) (-1.044,-1.632) (-1.064,-1.603) (-1.083,-1.574) (-1.101,-1.544) (-1.119,-1.513) (-1.136,-1.483) (-1.153,-1.452) (-1.169,-1.421) (-1.185,-1.389) (-1.200,-1.358) (-1.215,-1.326) (-1.229,-1.294) (-1.243,-1.262) (-1.256,-1.229) (-1.269,-1.196) (-1.282,-1.164) (-1.294,-1.131) (-1.306,-1.098) (-1.317,-1.064) (-1.328,-1.031) (-1.338,-0.997) (-1.348,-0.964) (-1.358,-0.930) (-1.368,-0.896) (-1.377,-0.862) (-1.386,-0.828) (-1.395,-0.794) (-1.403,-0.760) (-1.411,-0.726) (-1.420,-0.692) (-1.428,-0.658) (-1.436,-0.624) (-1.445,-0.590) (-1.454,-0.556) (-1.464,-0.522) (-1.462,-0.493) (-1.431,-0.478) (-1.398,-0.465) (-1.365,-0.454) (-1.331,-0.445) (-1.297,-0.437) (-1.262,-0.431) (-1.228,-0.425) (-1.193,-0.420) (-1.158,-0.416) (-1.123,-0.412) (-1.088,-0.409) (-1.053,-0.406) (-1.018,-0.404) (-0.983,-0.402) (-0.948,-0.400) (-0.913,-0.398) (-0.878,-0.396) (-0.843,-0.395) (-0.808,-0.394) (-0.772,-0.393) (-0.737,-0.392) (-0.702,-0.391) (-0.667,-0.390) (-0.632,-0.389) (-0.597,-0.388) (-0.562,-0.388) (-0.527,-0.387) (-0.492,-0.387) (-0.456,-0.386) (-0.421,-0.386) (-0.386,-0.386) (-0.351,-0.385) (-0.316,-0.385) (-0.281,-0.385) (-0.246,-0.385) (-0.211,-0.384) (-0.176,-0.384) (-0.140,-0.384) (-0.105,-0.384) (-0.070,-0.384) (-0.035,-0.384) (0.000,-0.384) (0.035,-0.384) (0.070,-0.384) (0.105,-0.384) (0.140,-0.384) (0.176,-0.384) (0.211,-0.384) (0.246,-0.385) (0.281,-0.385) (0.316,-0.385) (0.351,-0.385) (0.386,-0.386) (0.421,-0.386) (0.456,-0.386) (0.492,-0.387) (0.527,-0.387) (0.562,-0.388) (0.597,-0.388) (0.632,-0.389) (0.667,-0.390) (0.702,-0.391) (0.737,-0.392) (0.772,-0.393) (0.808,-0.394) (0.843,-0.395) (0.878,-0.396) (0.913,-0.398) (0.948,-0.400) (0.983,-0.402) (1.018,-0.404) (1.053,-0.406) (1.088,-0.409) (1.123,-0.412) (1.158,-0.416) (1.193,-0.420) (1.228,-0.425) (1.262,-0.431) (1.297,-0.437) (1.331,-0.445) (1.365,-0.454) (1.398,-0.465) (1.431,-0.478) (1.462,-0.493) (1.464,-0.522) (1.454,-0.556) (1.445,-0.590) (1.436,-0.624) (1.428,-0.658) (1.420,-0.692) (1.411,-0.726) (1.403,-0.760) (1.395,-0.794) (1.386,-0.828) (1.377,-0.862) (1.368,-0.896) (1.358,-0.930) (1.348,-0.964) (1.338,-0.997) (1.328,-1.031) (1.317,-1.064) (1.306,-1.098) (1.294,-1.131) (1.282,-1.164) (1.269,-1.196) (1.256,-1.229) (1.243,-1.262) (1.229,-1.294) (1.215,-1.326) (1.200,-1.358) (1.185,-1.389) (1.169,-1.421) (1.153,-1.452) (1.136,-1.483) (1.119,-1.513) (1.101,-1.544) (1.083,-1.574) (1.064,-1.603) (1.044,-1.632) (1.024,-1.661) (1.004,-1.690) (0.983,-1.718) (0.961,-1.745) (0.939,-1.772) (0.916,-1.799) (0.892,-1.825) (0.868,-1.851) (0.843,-1.876) (0.818,-1.900)};
\addplot[lineaD, flechaen=0.500] coordinates {(-0.576,-0.465) (-0.541,-0.464) (-0.506,-0.463) (-0.470,-0.463) (-0.435,-0.462) (-0.400,-0.462) (-0.365,-0.462) (-0.330,-0.461) (-0.294,-0.461) (-0.259,-0.461) (-0.224,-0.461) (-0.189,-0.460) (-0.154,-0.460) (-0.118,-0.460) (-0.083,-0.460) (-0.048,-0.460) (-0.013,-0.460) (0.022,-0.460) (0.058,-0.460) (0.093,-0.460) (0.128,-0.460) (0.163,-0.460) (0.198,-0.460) (0.234,-0.461) (0.269,-0.461) (0.304,-0.461) (0.339,-0.461) (0.374,-0.462) (0.410,-0.462) (0.445,-0.463) (0.480,-0.463) (0.515,-0.464) (0.550,-0.464) (0.586,-0.465) (0.621,-0.465) (0.656,-0.466) (0.691,-0.467) (0.726,-0.468) (0.762,-0.469) (0.797,-0.470) (0.832,-0.472) (0.867,-0.473) (0.902,-0.475) (0.937,-0.477) (0.973,-0.479) (1.008,-0.481) (1.043,-0.483) (1.078,-0.486) (1.113,-0.490) (1.148,-0.493) (1.182,-0.498) (1.183,-0.517) (1.156,-0.540) (1.130,-0.564) (1.104,-0.587) (1.077,-0.611) (1.051,-0.634) (1.025,-0.657) (0.998,-0.680) (0.971,-0.703) (0.945,-0.726) (0.918,-0.749) (0.890,-0.771) (0.863,-0.793) (0.835,-0.814) (0.807,-0.835) (0.778,-0.856) (0.749,-0.876) (0.720,-0.896) (0.691,-0.915) (0.661,-0.934) (0.630,-0.952) (0.600,-0.969) (0.569,-0.986) (0.537,-1.002) (0.506,-1.017) (0.474,-1.031) (0.441,-1.045) (0.408,-1.058) (0.375,-1.069) (0.342,-1.080) (0.308,-1.090) (0.274,-1.099) (0.240,-1.107) (0.205,-1.114) (0.170,-1.120) (0.136,-1.125) (0.101,-1.129) (0.066,-1.131) (0.030,-1.133) (-0.005,-1.133) (-0.040,-1.133) (-0.075,-1.131) (-0.110,-1.128) (-0.145,-1.124) (-0.180,-1.119) (-0.215,-1.113) (-0.249,-1.105) (-0.283,-1.097) (-0.317,-1.088) (-0.351,-1.077) (-0.384,-1.066) (-0.417,-1.054) (-0.450,-1.041) (-0.482,-1.027) (-0.514,-1.013) (-0.546,-0.997) (-0.577,-0.981) (-0.608,-0.964) (-0.639,-0.947) (-0.669,-0.928) (-0.699,-0.910) (-0.728,-0.890) (-0.757,-0.871) (-0.786,-0.850) (-0.814,-0.829) (-0.842,-0.808) (-0.870,-0.787) (-0.898,-0.765) (-0.925,-0.742) (-0.952,-0.720) (-0.979,-0.697) (-1.005,-0.674) (-1.032,-0.651) (-1.058,-0.628) (-1.085,-0.604) (-1.111,-0.581) (-1.137,-0.557) (-1.164,-0.534) (-1.190,-0.511) (-1.173,-0.496) (-1.138,-0.492) (-1.103,-0.489) (-1.068,-0.486) (-1.033,-0.483) (-0.998,-0.480) (-0.963,-0.478) (-0.928,-0.476) (-0.893,-0.474) (-0.857,-0.473) (-0.822,-0.471) (-0.787,-0.470) (-0.752,-0.469) (-0.717,-0.468) (-0.681,-0.467) (-0.646,-0.466) (-0.611,-0.465) (-0.576,-0.465)};
\addplot[lineaD, flechaen=0.164, flechaen=0.836] coordinates {(-3.200,0.000) (3.200,0.000)};

  \draw[texto, line width=0.6pt] (axis cs:-1.5,-0.5) rectangle (axis cs:1.5,0.5);
\end{axis}
\end{tikzpicture}
\caption{$\vD(\vr)$}
\end{subfigure}
\caption{Electrete cilíndrico uniformemente polarizado ($a/\ell=1/3$), calculado con las funciones de
corriente exactas; trazos: equipotenciales. (a) La polarización es uniforme y axial. (b) $\vE$ nace en
la cara positiva y muere en la negativa, de modo que dentro se opone a $\vP$; fuera es el campo de dos
discos con cargas $\pm\sigmaP$. (c) $\vD$ forma líneas cerradas, como el campo de un solenoide: dentro tiene
el sentido de $\vP$ y se curva junto a la pared lateral, donde $\rot\vP\neq\vcero$.}
\label{fig:cilindro-lineas}
\end{figure}
